% Source-derived chapter 05: Global Hecke action and closure relations
% Original source: intersection-complexes-unramified-l-factors
\chapter{Global Hecke action and closure relations}
\label{intersection-complexes-unramified-l-factors-chapter-05}
\source{intersection-complexes-unramified-l-factors}

\begin{remark}[Source section]
\label{intersection-complexes-unramified-l-factors-chapter-05-source}
\source{intersection-complexes-unramified-l-factors}
\source{intersection-complexes-unramified-l-factors}
This chapter reproduces the corresponding section of the retrieved
LaTeX source, including its definitions, statements, examples, and
proofs. It has not yet been translated into Lean.
\notready
\end{remark}

% The source section title was: Global Hecke action and closure relations
\label{sect:Heckeact}
\source{intersection-complexes-unramified-l-factors}


\emph{For the rest of this paper, assume that $\ch G_X = \ch G$
and all simple roots of $G$ are spherical roots of type $T$.} Equivalently, 
we are assuming that $B$ acts simply transitively on $X^\circ$ and 
for every simple root $\alpha$ of $G$, the $\PGL_2$-variety $X^\circ P_\alpha/\mathfrak R(P_\alpha)$ is isomorphic to $\mathbb G_m\backslash\PGL_2$ (over the algebraically closed field $k$).

\smallskip

As a consequence, $\Cal V \cap \ch\Lambda_X = \ch\Lambda^-_G$, the monoid of antidominant coweights of $G$. 
Recall from \S\ref{subsection:typeT} that 
the type $T$ assumption also implies that for every simple root $\alpha$, 
the open $P_\alpha$-orbit 
$X^\circ P_\alpha$ is the union of $X^\circ$ and the open $B$-orbits of 
two colors $\Cal D(\alpha) = \{ D_\alpha^+,D_\alpha^- \}$. 
We will let $\ch\nu_\alpha^\pm$ denote the valuation of $D_\alpha^\pm$, respectively.
Then $\ch\nu_\alpha^+ + \ch\nu_\alpha^- = \ch\alpha$ 
and $\brac{\alpha, \ch\nu_\alpha^\pm} = 1$. 
We encourage the reader to refer to Examples~\ref{eg:YHecke}, \ref{eg:YHeckePGL}, \ref{eg:Heckecompact} and \ref{eg:Yfiber}.

\subsection{Main results of this section}

This section is quite technical, and we advise the reader to read the main results listed here, and skip the rest of the section at first reading. Before we introduce the results, let us observe that, so far, we have uniformly treated all affine spherical varieties. 
However, the classification of spherical varieties is divided into two parts: 
(i) the classification of homogeneous spherical varieties $H\bs G$, and (ii) the
classification of spherical embeddings $H\bs G \into X$ (by a spherical embedding we mean a $G$-equivariant open, dense embedding $H\bs G \into X$, where $X$ is a normal, and hence spherical, $G$-variety). 

Since $X$ is affine, $X^\bullet=H\bs G$ is quasiaffine and there is a 
\emph{canonical affine embedding} 
\[ X^\can := \Spec k[H\bs G]. \] 
(The fact that the coordinate ring of $k[H\bs G]$ is finitely generated follows from the fact that $B$-eigenspaces are one-dimensional, and the $B$-character group is finitely generated.) By normality, $X^\can$ has no divisors that do not meet $H\bs G$, 
i.e., $\Cal D(X)$ has no $G$-stable divisors, and coincides with the set $\Cal D$ of colors. 
Therefore, the cone $\Cal C_0(H\bs G):= \Cal C_0(X^\can)$ 
is generated by the valuations $\varrho_X(\Cal D)$ of colors.

For any other affine spherical embedding $H\bs G\into X$, there is  
a natural map $X^\can \to X$, so $X^\can$ is universal among affine embeddings 
of $H\bs G$. 

It turns out that this distinction between the minimal and the general embeddings is important when we consider arc spaces and their global models. 
As we will recall in Lemma~\ref{lem:Membedding}, the map $X^\can \to X$ induces a closed embedding of mapping stacks 
\[ \Scr M_{X^\can} \hookrightarrow \Scr M_X,\]
whose image is a union of irreducible components.

We are more interested in the closure of $\Scr M^0_{X^\bullet} = \Bun_H$ in the former, which we will denote by $\barM_X^0$ (it may or may not be the same as $\Scr M_{X^\can}$ depending on whether the monoid $\mathfrak c_{X^\can}$ is generated by colors, see Corollary \ref{cor:can-comp}):
\begin{equation}
 \label{cantogeneral}
\source{intersection-complexes-unramified-l-factors}
\barM_X^0 \hookrightarrow \Scr M_{X^\can} \hookrightarrow \Scr M_X.
\end{equation}



The main theme of this section is, in some sense, a reconstruction of $\Scr M_X$ out of suitable Hecke operators acting on $\barM_X^0$: 
For any $\ch\Theta \in \Sym^\infty(\mathfrak c_X^-\sm 0)$, a multiset of nonzero elements in $\mathfrak c_X^-$,
there is a natural proper map 
\begin{equation} \label{equation-action}
\source{intersection-complexes-unramified-l-factors}
\act_\sM : \barM_X^0 \ttimes \ol\Gr^{\ch\Theta}_{G, C^{\ch\Theta}} \to 
\Scr M_X,
\end{equation}
which corresponds to the action on the ``basic stratum'' $\barM_X^0$ by the closed stratum of the affine Grassmannian parametrized by $\ch\Theta$.
(See Proposition--Construction \ref{prop-const:action}.) Recall that such a multiset $\check\Theta$ also parametrizes a stratum $\Scr M^{\ch\Theta}_X$ of the global mapping space (\S \ref{sect:globstrat}). The main technical result of this section is the following:

\begin{thm}
\source{intersection-complexes-unramified-l-factors}
\notready
\label{thm:comp}
\source{intersection-complexes-unramified-l-factors}
For every $\ch\Theta \in \Sym^\infty(\mathfrak c_X^- \sm 0)$, the action map \eqref{equation-action} 
\begin{enumerate}
\item has image equal to the closure of $\Scr M^{\ch\Theta}_X$, and  
\item is birational onto its image. 
\item The restriction of $\act_\sM$ to $\act_\sM^{-1}(\sM^{\ch\Theta}_X) \cap 
(\Bun_H \ttimes \Gr^{\ch\Theta}_{G, \oo C^{\ch\Theta}}) \to \sM^{\ch\Theta}_X$ 
is an isomorphism. 
\end{enumerate}
\end{thm}

The proof of Theorem \ref{thm:comp} will be given in \S \ref{proof:comp}.

We use this theorem to achieve two goals in this section: 


The first goal is to understand irreducible components of $\Scr M_X$ and closure relations among the strata $\Scr M^{\ch\Theta}_X$. 
We will introduce the natural generalization to multisets of the order $\succeq$ among elements of the lattice $\ch\Lambda_X$ (we remind that this order is determined by the monoid of colors, see \S \ref{sect:spherical}), and prove:
\begin{prop}[See Proposition \ref{prop:closure-rel}.]
\source{intersection-complexes-unramified-l-factors}
\notready
\label{proposition-closures}
\source{intersection-complexes-unramified-l-factors}
Let $\ch\Theta,\ch\Theta' \in \Sym^\infty(\mathfrak c_X^-\sm 0)$. 
The stratum $\sM^{\ch\Theta'}_X$ lies in the closure of $\sM^{\ch\Theta}_X$
if and only if there exists $\ch\Theta''$  such that $\ch\Theta$ refines $\ch\Theta''$ and $\ch\Theta' \succeq \ch\Theta''$. 
\end{prop}

Moreover, observe that the irreducible components of $\barM_X^0$ are in bijection with connected components of $\Scr M_{X^\bullet} = \Bun_H$, i.e., parametrized by $\pi_0(\Bun_H) = \pi_1(H)$. Using the action \eqref{equation-action} on those components gives us a parametrization of the irreducible components of the closure of each stratum:
\begin{prop}[See Corollaries \ref{cor:strata-comp} and \ref{cor:Ycomp}]
\source{intersection-complexes-unramified-l-factors}
\notready
\label{proposition-irred-strata}
\source{intersection-complexes-unramified-l-factors}
For every $\ch\Theta \in \Sym^\infty(\mathfrak c_X^-\sm 0)$, the irreducible components in the closure $\barM_X^{\ch\Theta}$ of the corresponding stratum are naturally parametrized by $\pi_1(H)$. 
For any $\ch\lambda \in \mf c_X$, the base change of 
an irreducible component to $\sY^{\ch\lambda} \xt_{\sM_X} 
\barM_X^{\ch\Theta}$ is still irreducible (when nonempty).
\end{prop}

The combination of Propositions \ref{proposition-closures} and \ref{proposition-irred-strata} implies:

\begin{cor}[See Corollary \ref{cor:comp}]
\source{intersection-complexes-unramified-l-factors}
\notready 
There is a natural bijection between the set of irreducible
components of $\Scr M_X$ and 
\[ \pi_1(H) \xt \Sym^\infty(\Cal D^G_{\mathrm{sat}}(X)),\]
where $\Cal D^G_{\mathrm{sat}}(X)$ denotes the set of primitive elements in $\mathfrak c_X^-$ that
cannot be decomposed as a sum $\ch\theta+\ch\nu_D$ where 
$\ch\theta\in  \mathfrak c_X^-\sm 0$ and $\ch\nu_D$ is the valuation attached to a color. 
\end{cor}



The second goal achieved by Theorem \ref{thm:comp} is to reduce the study of the IC complex of an arbitrary mapping space $\Scr M_X$ to that of the minimal affine embedding:

\begin{thm}
\source{intersection-complexes-unramified-l-factors}
\notready \label{thm:heckeIC}
\source{intersection-complexes-unramified-l-factors}
For any $\ch\Theta \in \Sym^\infty(\mathfrak c_X^- \sm 0)$, there
is a natural isomorphism 
\[ \IC_{\barM_X^0} \star \IC_{\ol\Gr^{\ch\Theta}_{G,C^{\ch\Theta}}} \cong 
\IC_{\barM^{\ch\Theta}_X}. \] 
\end{thm}

We point to \S\ref{sect:pf:heckeIC} for the notation and the proof.  
Theorem~\ref{thm:heckeIC} and its proof are independent from the rest of this paper; we
include it only for conceptual completeness. 
The proof involves a passage to Zastava model (the extra structure of a flag)
and uses the results of the sequel \S\ref{sect:centralfiber}.
The argument of proof can also be extended to meromorphic quasimaps to verify
\cite[Conjecture 7.3.2]{GN} in the case $\ch G_X = \ch G$. 

\subsection{Hecke action on global model} \label{sect:Heckeglob}
\source{intersection-complexes-unramified-l-factors}

Now we introduce the action map \eqref{equation-action}, as an analog of the 
action of $G(F)$ on $X^\bullet(F)$ for the global model.

The notation is cumbersome because we give a multi-point version of the action, but
the idea is simple: in the notation of \S\ref{sect:adelic}, 
the $k$-points of $\sM_X$ correspond to a subset of $H(\Bbbk) \bs G(\mbb A) / G(\mbb O)$, where $\mbb O = \prod_{v\in \abs C} \mf o_v$. 
We have a Hecke correspondence 
\[ H(\Bbbk) \bs G(\mbb A) \xt^{G(\mbb O)} G(\mbb A) / G(\mbb O) \to H(\Bbbk) \bs G(\mbb A) / G(\mbb O) \]
induced by multiplication in $G$. We think of $G(\mbb A)/G(\mbb O)$ as a
factorizable version of the affine Grassmannian. Then the ``action map'' 
is simply the restriction of the above to a positively graded subset of $G(\mbb A)/G(\mbb O)$ 
such that everything maps to $(X^\bullet(\mbb A) \cap X(\mbb O))/G(\mbb O)$, i.e., 
the locus of regular maps. 

\subsubsection{Positively graded subscheme of factorizable affine Grassmannian}\label{sect:strata-Grassmannian}
\source{intersection-complexes-unramified-l-factors}
Let $\ch\Theta = \sum_{\ch\theta} N_{\ch\theta} 
[\ch\theta] \in \Sym^\infty(\mathfrak c_X^- \sm 0)$ 
be a partition. We have a closed 
subscheme
\[ \ol\Gr^{\ch\theta}_{G,C} = \ol\Gr_G^{\ch\theta} \ttimes C := \ol\Gr_G^{\ch\theta} \xt^{\on{Aut}^0( k\tbrac t )} \on{Coord}^0(C) \subset \Gr_{G,C} \] 
where $\Aut^0(k\tbrac t) = \Spec k[a_1^{\pm 1}, a_2,\dotsc]$ is the group scheme of algebra automorphisms of $k\tbrac t$ that preserve the maximal ideal, and $\on{Coord}^0(C)\to C$ is the $\Aut^0(k\tbrac t)$-torsor classifying $v\in C$ together with an isomorphism $k\tbrac t \cong \mf o_v$ sending $t$ to a uniformizer 
(see \S\ref{sect:Llocalization} or \cite[(3.1.11)]{Xinwen}).
Consider the $N_{\ch\theta}$-fold product $(\ol\Gr_{G,C}^{\ch\theta})^{N_{\ch\theta}} \xt_{C^{N_{\ch\theta}}} {\oo C^{N_{\ch\theta}}}$ restricted to the disjoint locus with all diagonals removed. 
This descends to a subscheme $\ol\Gr^{N_{\ch\theta}\ch\theta}_{G, \oo C^{(N_{\ch\theta})}} \subset \Gr_{G, \oo C^{(N_{\ch\theta})}}$. 
In the notation of \S\ref{def:partition}, let 
\[ \ol\Gr^{\ch\Theta}_{G, \oo C^{\ch\Theta}} := \oo\prod_{\ch\theta} \ol\Gr^{N_{\ch\theta}\ch\theta}_{G, \oo C^{(N_{\ch\theta})}} 
\subset \Gr_{G, C^{(\abs{\ch\Theta})}} \xt_{C^{(\abs{\ch\Theta})}} 
\oo C^{\ch\Theta} \]
and let $\ol\Gr^{\ch\Theta}_{G,C^{\ch\Theta}}$ denote its closure 
in $\Gr_{G,C^{(\abs{\ch\Theta})}} \xt_{C^{(\abs{\ch\Theta})}} C^{\ch\Theta}$.
We consider $\ol\Gr^{\ch\Theta}_{G, C^{\ch\Theta}}$ as a scheme 
over $C^{(\abs{\ch\Theta})}$ with an action of the group scheme 
$(\Scr L^+ G)_{C^{(\abs{\ch\Theta})}}$, the multi-point version of the arc space 
defined in \S\ref{sect:multijet}. 

\medskip

The closure relations of the Beilinson--Drinfeld affine Grassmannian 
are known (cf.~\cite[Proposition 3.1.14]{Xinwen}), so we can describe
the reduced fiber of $\ol\Gr^{\ch\Theta}_{G,C^{\ch\Theta}}$ 
over a point of $C^{\ch\Theta}$ as follows: 
a point of $C^{\ch\Theta}$ is the collection $(D^{\ch\theta})_{\ch\theta}$, 
for each $\ch\theta$,
of a degree $N_{\ch\theta}$ divisor 
$D^{\ch\theta} = \sum_{v\in \abs C} N_{\ch\theta,v} v$. 
The reduced fiber of $\ol\Gr^{\ch\Theta}_{G,C^{\ch\Theta}}$ 
over this point is the scheme
\begin{equation} \label{e:GrTheta-fiber}
\source{intersection-complexes-unramified-l-factors}
\prod_{v\in \abs C} \ol\Gr^{\sum_{\ch\theta} N_{\ch\theta,v}\ch\theta}_{G,v}. 
\end{equation}


\subsubsection{}  \label{sect:Glevel}
\source{intersection-complexes-unramified-l-factors}
Let $\widehat\sM_X$ denote the stack 
representing the data of 
\[ (\sigma,\Scr P_G)\in \Scr M_X,\, D\in \Sym C,
\text{ and a trivialization }\Scr P_G|_{\wh C'_D} \cong \Scr P^0_G|_{\wh C'_D}
\] 
(see \S\ref{sect:BLthm} for the definition of $\wh C'_D$), i.e., 
a point of $\sM_X$ together with infinite $G$-level structure at points in the support of $D$. 
This admits a natural action by $\Scr L^+ G$, and 
the forgetful map $\widehat{\Scr M}_X
 \to \Scr M_X \xt \Sym C$ is a $\Scr L^+G$-torsor.
Let 
\[ \Scr M_X \ttimes \ol\Gr^{\ch\Theta}_{G, C^{\ch\Theta}} := 
\widehat{\Scr M}_X \xt^{\Scr L^+ G}_{\Sym C}
\ol\Gr^{\ch\Theta}_{G,C^{\ch\Theta}} \]
denote the twisted product over $C^{(\abs{\ch\Theta})} \subset \Sym C$.

For an affine test scheme $S$, 
an $S$-point of $\Scr M_X \ttimes \ol\Gr^{\ch\Theta}_{G,C^{\ch\Theta}}$ 
is the data 
$(\sigma, \Scr P_G,\Scr P'_G, (D^{\ch\theta})_{\ch\theta}, \tau)$
where 
\begin{itemize}
\item $\Scr P_G,\Scr P_G'$ are $G$-bundles on $C\xt S$,
\item $\sigma : C\xt S \to X\xt^G \Scr P_G$ is a section such that $(\sigma,\Scr P_G)\in \Scr M_X(S)$,
\item for each $\ch\theta \in \mathfrak c_X^- \sm 0$,  
we have a degree $N_{\ch\theta}$ effective Cartier divisor 
$D^{\ch\theta} \subset C\xt S$, 
\item $\tau : \Scr P'_G|_{C\xt S\sm D} \cong \Scr P_G|_{C\xt S\sm D}$
is a trivialization for $D:=\sum_{\ch\theta} D^{\ch\theta}$
\end{itemize}
such that after fixing an isomorphism 
$\Scr P_G|_{\wh C'_D} \cong \Scr P^0_G|_{\wh C'_D}$ (which always exists
after flat base change over $S$), the datum
$((D^{\ch\theta}),\Scr P'_G|_{\wh C'_D},\tau)$ defines
a point in $\ol\Gr^{\ch\Theta}_{G,C^{\ch\Theta}}$.
Here we are implicitly using Beauville--Lazlo's theorem to pass between
the global and local descriptions of $\Gr_{G,\Sym C}$, 
cf.~\S\ref{sect:localGr}.


\begin{prop-const}
\label{prop-const:action}
\source{intersection-complexes-unramified-l-factors}
For any $\ch\Theta \in \Sym^\infty(\mathfrak c_X^-\sm 0)$
there is a natural proper map 
\[ \act_\sM : \Scr M_X \ttimes \ol\Gr^{\ch\Theta}_{G, C^{\ch\Theta}} \to 
\Scr M_X. \]
\end{prop-const}

\begin{proof}
Fix an affine test scheme $S$ and $(\sigma,\Scr P_G,\Scr P'_G, (D^{\ch\theta})_{\ch\theta},\tau) \in \Scr M_X \ttimes \ol\Gr^{\ch\Theta}_{G,C^{\ch\Theta}}(S)$. 
The composition $\tau^{-1} \circ \sigma|_{C\xt S\sm D}$ defines
a section $\sigma' : C\xt S \sm D \to X\xt^G \Scr P'_G$. 
We claim that $\sigma'$ extends to a regular map on $C\xt S$. 
Given this claim, we can define $(\sigma',\Scr P'_G)\in \Scr M_X(S)$ 
to be the image of $\act_\sM$. Properness of $\act_\sM$
follows from properness of $\ol\Gr^{\ch\Theta}_{G,C^{\ch\Theta}}$
and Lemma~\ref{lem:extendsection}.

We now prove the claim that $\sigma'$ extends to all of $C\xt S$.
Since $k[X]$ is a locally finite $G$-module, it is generated as an
algebra by some finite dimensional $G$-submodule $V \subset k[X]$. 
This induces a $G$-equivariant embedding of varieties $\varphi: X\into V^*$,
where $V^{*}$ is considered as a right $G$-module. 
It suffices to show that the composition $\varphi(\sigma'):
C\xt S\sm D \to V^{*} \xt^G \Scr P'_G$ extends. 
For any weight $\mu$ of $V \subset k[X]$, we have $\mu \le \lambda$ for $\lambda\in \mf c_X^\vee$. 
Therefore given $\ch\theta\in \mathfrak c_X^-$, we have $\brac{\mu,\ch\theta}\ge 0$
for all weights $\mu$ of $V$. 
We deduce that the group homomorphism $G \to \GL(V)$ induces a natural map
\[ \ol\Gr^{\ch\Theta}_{G,C^{\ch\Theta}} \to 
\Scr L^+\mathrm{End}(V)/\Scr L^+\GL(V). \]
Hence $\tau^{-1}$ induces a regular map 
$V^{*} \xt^G\Scr P_G|_{\wh C'_D} \to V^{*} \xt^G\Scr P'_G|_{\wh C'_D}$
and $\varphi(\tau^{-1}\circ \sigma|_{\wh C'_D})$ defines a section 
$\wh C'_D \to V^{*} \xt^G \Scr P'_G$. 
By Beauville--Laszlo's theorem (cf.~\cite[Theorem 2.12.1]{BD}), 
this implies that $\varphi(\sigma')$ is defined on all of $C\xt S$.
\end{proof}


We describe more precisely what $\act_\sM$ is doing on $k$-points: 
at a single $v \in \abs C$, a $k$-point of $\Scr M_X$ gives an element
of $(X(\mf o_v)\cap X^\bullet(F_v)) / G(\mf o_v)$ and the map $\act_\sM$ corresponds
to the natural $G(F_v)$-action on $X^\bullet(F_v)$.
For $\ch\mu \in \ch\Lambda^-_G$, 
define the set  
\[ X^\bullet(F_v)_{G:\succeq \ch\mu} = \bigcup_{\substack{\ch\mu' \in \ch\Lambda^-_G, \ch\mu' \succeq \ch\mu}} X^\bullet(F_v)_{G:\ch\mu'}, \]
in the notation of \S \ref{sect:stratLXpoint}.
In the next subsection we prove a slightly more precise\footnote{In \cite{SV} 
the ordering $\succeq$ is defined with respect to the \emph{rational} cone generated by
the valuations of colors, whereas we define $\succeq$ with respect
to the monoid generated by non-negative \emph{integral} combinations of valuations of colors.} 
version of \cite[Lemma 5.5.2]{SV}:

\begin{lem}
\source{intersection-complexes-unramified-l-factors}
\notready \label{lem:act-kpoints}
\source{intersection-complexes-unramified-l-factors}
Let $\ch\mu,\ch\theta \in \ch\Lambda^-_G$. 
The action map sends 
\[
    X^\bullet(F_v)_{G:\succeq \ch\mu} \xt^{G(\mf o_v)} \ol{\msf L^+G \cdot t^{\ch\theta} \cdot \msf L^+G}(k) \to X^\bullet(F_v)_{G:\succeq \ch\mu+ \ch\theta}. 
\]
\end{lem}


\subsection{Reduction to $\mf c_{X^\bullet} = \mbb N^{\Cal D}$} \label{sect:freemonoid}
\source{intersection-complexes-unramified-l-factors}
We are interested in applying Hecke actions to $\barM_X^0$, the closure of $\sM_{X^\bullet}=\sM_X^0=\Bun_H$ in $\sM_X$, since it is the most basic closure of a stratum. 
On the other hand, in order to determine the stratification of $\barM_X^0$ we
need a moduli description of this stack. A first guess would be 
that $\barM_X^0 = (\sM_{X^\can})_\red$, but this may not be true if $\mf c_{X^\bullet} := \mf c_{X^\can}$ 
is not equal to $\mathfrak c_X^{\Cal D} = \{\ch\lambda\succeq 0\}$.
In this subsection we explain how to get around this technical issue: 
we can always replace $G$ by a central extension $G'\to G$ such that if
$H\bs G = H'\bs G'$ and $X' := \Spec k[H'\bs G']$, then 
$\barM_{X'}^0 = \sM_{X'}$ and $\mf c_{X'} = \mbb N^{\Cal D}$ (see Lemma~\ref{lem:Xcover}).

\smallskip 

This is a generalization of the need to replace $\ol{N\bs G}^\aff$ by 
$\ol{N\bs \wt G}^\aff$ for $\wt G$ a simply connected cover of $G$ 
in \cite[\S 4.1]{ABB}, \cite[\S 7.2]{Sch} 
to correctly define Drinfeld's compactification of $\Bun_N$ 
for an arbitrary reductive group $G$.

\subsubsection{} 

Let us first assume that $k[G]$ is a UFD. Further, assume that $H$ is connected, as is the case under our assumptions (Remark~\ref{rem:Hconnected}).
Then the preimage of any color $D$ in $G$ is irreducible, and this defines 
a bijection between $H\times B$-stable prime divisors in $G$ and colors; moreover, under our UFD assumption, the former are all principal. 

Let $k(G)^{(H\xt B)}$ denote the $H\xt B$-eigenvectors 
of $k(G)$. 
Then, since $H B$ is open dense in $G$, we have
\[ k(G)^{(H\times B)}/k^\times = \Cal X(H) \xt_{\Cal X(B\cap H)} \Cal X(B), 
\] 
where $\Cal X(H)$ is the character group of $H$ 
(so $\Cal X(B)=\Lambda_{G}$ by definition). The valuation map gives rise to a short exact sequence
\begin{equation} \label{e:color-basis}
\source{intersection-complexes-unramified-l-factors}
 0 \to \Cal X(G) \to  
    \Cal X(H) \xt_{\Cal X(B \cap H)} \Cal X(B) \to   \mbb Z^{\Cal D} \to 0,
\end{equation}
which sends an element of $k(G)^{(H\times B)}/k^\times$ to its divisor, identified with a $\mathbb Z$-linear combination of the colors. (We have used here both the UFD property, and the fact that invertible regular functions on $G$ are multiples of characters.)


The preimage in $G$ of each color $D$ defines a valuation $v_{HD}$ on $k(G)^\times$; its restriction to $k(G)^{(H\xt B)}$ defines a map
\[ 
    \wt\varrho_{H\bs G} : \Cal D \to (\Cal X(H) \xt_{\Cal X(B\cap H)} \Cal X(B))^\vee. 
\]
Composing with the second projection to $\ch\Lambda_X$ gives the usual
valuation map $\varrho_{H\bs G}$.
By construction we have $v_{HD}(f_{D'}) = \delta_{D,D'}$ 
for $D,D'\in \Cal D$, and the sequence dual to \eqref{e:color-basis} is 
\[ 
    0\to \mbb Z^{\Cal D} \overset{\wt\varrho_{H\bs G}}\longrightarrow 
(\Cal X(H) \xt_{\Cal X(B\cap H)} \Cal X(B))^\vee \to 
\Cal X(G)^\vee \to 0.
\]

\subsubsection{} \label{sect:Xcover}
\source{intersection-complexes-unramified-l-factors}
Now we return to the case of arbitrary $G$. Let 
\[ 1 \to Z \to \wt G \to G \to 1 \] 
be a central extension with connected 
kernel $Z$ such that the derived group $[\wt G, \wt G]$ is simply
connected (such a $\wt G$ always exists). 
Then $k[\wt G]$ is a UFD (\cite[Proposition 4.6]{KKLV}, \cite{Iv}).
We consider $H\bs G$ as a $\wt G$-spherical variety, so 
$H\bs G = \wt H \bs \wt G$ where $\wt H = HZ$ is the preimage of $H$
in $\wt G$. Let $\wt B$ denote the Borel subgroup of $\wt G$.
If $H$ is connected, then $\wt H$ is also connected.
The colors stay the same, so 
$\wt\varrho_{\wt H\bs \wt G}$ induces a short exact sequence
\begin{equation} \label{e:color-ses} 
\source{intersection-complexes-unramified-l-factors}
0 \to \mbb Z^{\Cal D} \overset{\wt\varrho_{\wt H\bs \wt G}}\longrightarrow 
(\Cal X(\wt H) \xt_{\Cal X(\wt B\cap \wt H)} \Cal X(\wt B))^\vee \to 
\Cal X(\wt G)^\vee \to 0.
\end{equation}

\smallskip

Let $\wt H_{\mathrm{ab}} = \wt H/[\wt H,\wt H]$, so 
$\Cal X(\wt H) = \Cal X(\wt H_{\mathrm{ab}})$. 
We consider $[\wt H,\wt H] \bs \wt G$ 
as a spherical variety for the group $G' := \wt H_{\mathrm{ab}} \overset Z\times \wt G$,
where $\wt H_{\mathrm{ab}}$ acts by \emph{left} translation and 
$\wt G$ acts by right translation. 
Then 
$[\wt H,\wt H] \bs \wt G = H' \bs G'$ 
where $H'=(\wt H/Z)^\diag=H^\diag$ is diagonally embedded in $G'$. 
Observe that $G' \to G$ is a central extension
with kernel $\wt H_{\mathrm{ab}}$
and 
\[ H'\bs G' = [\wt H,\wt H]\bs \wt G \to \wt H \bs \wt G = H\bs G \]
is a $\wt H_{\mathrm{ab}}$-torsor. 
The Borel subgroup $B'$ of $G'$ is $\wt H_{\mathrm{ab}} \xt^Z \wt B$
and $B'\cap H' = ((\wt B \cap \wt H)/Z)^\diag$, so 
 \eqref{e:color-ses} is equivalent to a short exact sequence
\[ 
    0 \to \mbb Z^{\Cal D} \overset{\varrho_{H'\bs G'}}\longrightarrow
    \ch\Lambda_{H'\bs G'} = \ker(\Cal X(B')\to \Cal X(B'\cap H'))^\vee 
    \to \Cal X(\wt G)^\vee \to 0. 
\]

To summarize:
\begin{lem}
\source{intersection-complexes-unramified-l-factors}
\notready \label{lem:Xcover}
\source{intersection-complexes-unramified-l-factors}
Let $H\bs G$ be a homogeneous spherical variety with $H$ connected. 
Then there exists a central extension $G' \to G$ 
with connected kernel $Z'$ and a spherical
subgroup $H'\subset G'$ such that 
\begin{enumerate}
\item $[G',G']$ is simply connected. 
\item The covering $G'\to G$ restricts to an isomorphism $H'\cong H$. 
In particular, 
the projection $H' \bs G' \to H \bs G$ is a $Z'$-torsor.
\item 
The valuation map $\varrho_{H'\bs G'}$ embeds 
$\mbb Z^{\Cal D} \into \ch\Lambda_{H'\bs G'}$
as a direct summand. 
\end{enumerate}
\end{lem}



\noindent By (iii), we have that $\mf c_{H'\bs G'} = \mathfrak c^{\Cal D}_{H'\bs G'}= \mbb N^{\Cal D}$ so there is no question of integrality. For a similar result, see \cite[Lemma 2.1.1]{Brion07}.

\begin{eg}
\source{intersection-complexes-unramified-l-factors}
\notready
If we start with $H = \mbb G_m$ the torus inside $G = \PGL_2$, then
$G' = \mbb G_m \xt \GL_2$ and $H'= \mbb G_m$ where $\mbb G_m$ 
maps to $\GL_2$ by $\left( \begin{smallmatrix} * & 0 \\ 0 & 1 \end{smallmatrix} \right)$. So $H'\bs G' = \GL_2$, where $\mbb G_m$ acts by 
left translations and $\GL_2$ by right translations.
This is a $\mbb G_m$-torsor over $\mbb G_m \bs \GL_2$ and a 
$\mbb G_m \xt Z(\GL_2)$-torsor over $\mbb G_m \bs \PGL_2$.
\end{eg}



\begin{proof}[{Proof of Lemma~\ref{lem:act-kpoints}}]
Let $1\to Z'\to G'\to G\to 1$ be a central extension as in Lemma~\ref{lem:Xcover}.  
Since $\mf c_{H'\bs G'} = \mathfrak c^{\Cal D}_{H'\bs G'}$, 
there is no question of integrality and \cite[Lemma 5.5.2]{SV} 
applied to the $G'$-variety $H'\bs G'$ says that 
the action map sends 
\[ (H'\bs G')(F_v)_{\succeq \tilde\mu} \xt^{G'(\mf o_v)} G'(\mf o_v) t^{\tilde\theta} G'(\mf o_v) \to (H'\bs G')(F_v)_{\succeq \tilde\mu+\tilde\theta}, \]
where $\tilde\mu, \tilde\theta \in \ch\Lambda_{H'\bs G'}^-$. 
Since $H'\bs G' \to H\bs G = X^\bullet$ is a $Z'$-torsor, 
it is surjective on $F_v$-points (and corresponding orbits). 
Since $\mathfrak c^{\Cal D}_{H'\bs G'}$ maps to $\mathfrak c_X^{\Cal D}$, 
we deduce that the action map sends 
\[ X^\bullet(F_v)_{\succeq \ch\mu} \xt^{G(\mf o_v)} G(\mf o_v) t^{\ch\theta} G(\mf o_v)
\to X^\bullet(F_v)_{\succeq \ch\mu+\ch\theta} \]
for $\ch\mu,\ch\theta \in \ch\Lambda^-_G$.
The preimage of $\ol\Gr^{\ch\theta}_G(k)$ in $G(F_v)$ consists of the union of 
$G(\mf o_v) t^{\ch\theta'} G(\mf o_v)$ for $\ch\theta'\in \ch\Lambda^-_G,\,
\ch\theta' \ge \ch\theta$. 
Since $\ch\theta' \ge \ch\theta$ implies $\ch\theta' \succeq \ch\theta$, 
we deduce the claim. 
\end{proof}



\subsection{Open Zastava} \label{sect:Y0cc}
\source{intersection-complexes-unramified-l-factors}
Consider $\Scr Y^{?,0} = \sY_{X^\bullet}=\sY \xt_{\sM_X} \sM_X^0 = \Maps_\gen(C, H\bs G/B \supset \pt)$, which 
is an open subscheme of $\sY$. 
Since $\sM_X^0 \cong \Bun_H$ is smooth and $\sY_{X^\bullet}$ is smooth locally isomorphic to 
$\sM_X^0$ by Lemma~\ref{lem:localglobalyoga}, the scheme $\sY_{X^\bullet}$ is also smooth. 

The preimage of the connected component $\sA^{\ch\lambda}$ in $\sY_{X^\bullet}$ 
is by definition $\sY^{\ch\lambda,0}$. 
Let $G' \to G$ be a central extension as in Lemma~\ref{lem:Xcover}, so 
we have a torsor $H'\bs G' \to X^\bullet$. 
Let $X' = \Spec k[H'\bs G']$. 
Then $\mathfrak c_{X'} = \mathbb N^{\Cal D}$ and we have an isomorphism of stacks
$X'^\bullet / B' \cong  X^\bullet/ B$.
Hence, 
\begin{equation}\label{e:YXprime}
\source{intersection-complexes-unramified-l-factors}
    \sY_{X^\bullet} \cong \sY_{X'^\bullet} = \Maps_\gen( X'^\bullet/B' \supset \pt ),
\end{equation}
and the map $\sY_{X^\bullet} \to \sA$ factors through $\pi_{X'} : \sY_{X'} \to \sA_{X'}$.
The base $\sA_{X'}$ is a disjoint union of smooth partially symmetrized powers of the curve indexed by $\mathbb N^{\Cal D}$. 
 For $D = \sum_{\Cal D} n_{D'}\cdot D'\in \mathbb N^{\Cal D}$, we will denote by 
$\sY_{X^\bullet}^D$ the preimage of the corresponding component of $\sA_{X'}$. 
We define $\varrho_X(D) = \sum n_{D'}\ch\nu_{D'}$ and $\len(D) = \sum n_{D'}$. Hence, if $\ch\lambda = \varrho_X(D)$, then $\sY_{X^\bullet}^D \subset \sY_{X^\bullet}^{\ch\lambda}$. 

This implies: 

\begin{lem}
\source{intersection-complexes-unramified-l-factors}
\notready \label{lem:Ynonempty}
\source{intersection-complexes-unramified-l-factors}
The stratum $\Scr Y^{\ch\lambda,0}$ is nonempty only if $\ch\lambda \succeq 0$.
\end{lem}



Under our assumption that all simple roots of $G$ are spherical roots of type $T$, 
every color $D\in \Cal D$ belongs to $\Cal D(\alpha)$ for some simple root $\alpha$ of $G$
and $X^\circ P_\alpha/\mf R(P_\alpha) = \mbb G_m \bs \PGL_2$. 

\begin{lem}
\source{intersection-complexes-unramified-l-factors}
\notready \label{lem:Ycolor}
\source{intersection-complexes-unramified-l-factors}
If $\Cal D(\alpha) = \{ D_\alpha^+, D_\alpha^- \}$ for a simple root $\alpha$,
then for any $n^\pm \in \mbb N$, 
there is an isomorphism 
\[ \sY^{n^+ D_\alpha^+ + n^- D_\alpha^-}_{X^\bullet} = C^{(n^+)}\oo\xt C^{(n^-)}. \]
\end{lem}
\begin{proof} 
We may assume $\mf c_{X^\bullet} = \mbb N^{\Cal D}$. Then, 
$\sY^{n^+ D_\alpha^+ + n^- D_\alpha^-}_{X^\bullet}$
classifies maps from the curve to $X^\bullet/B$ which have zero valuation on every color other than those in $\Cal D(\alpha)$, and therefore is a subscheme of $\Maps_\gen(C, (X^\bullet-\bigcup_{D\in \Cal D \sm \Cal D(\alpha)} D)/B \supset \pt)$. 
Since $X^\circ P_\alpha = \mathbb G_m\backslash P_\alpha$ is affine, its complement is the union of the colors that it does not intersect. Hence 
\[ (X^\bullet-\bigcup_{D\in \Cal D \sm \Cal D(\alpha)}D)/B = (X^\circ P_\alpha)/B = \mathbb G_m\bs\PGL_2/B_{\PGL_2} = \mbb G_m \bs \mbb P^1.\]
By Example~\ref{eg:YHecke}, we see that $\Maps_\gen(C, \mbb G_m \bs \mbb P^1 \supset \pt) = \Sym C \oo\xt \Sym C$,
and it follows that the components correspond to $\mbb N^{\Cal D(\alpha)}$ in
the natural way.
\end{proof}



\begin{rem}
\source{intersection-complexes-unramified-l-factors}
\notready \label{rem:Ycomp}
\source{intersection-complexes-unramified-l-factors}
For a general element $D = \sum_{\Cal D} n_{D'}\cdot D' \in \mbb N^{\Cal D}$, 
the graded factorization property together with Lemma~\ref{lem:Ycolor} imply
that there is an open embedding 
$\oo\prod_{\Cal D} C^{(n_{D'})} \into \sY^D_{X^\bullet}$,
where the product is over the disjoint locus. 
We will show in Lemma~\ref{lem:oYconnected} that 
$\sY^D_{X^\bullet},\, D\in \mbb N^{\Cal D}$ are precisely the connected components
of $\sY_{X^\bullet}$, so the open subscheme above is dense. 
(This also implies that the $\sY^D_{X^\bullet}$ are defined intrinsically and independently
of the choice of $G'\to G$.)
\end{rem}

\begin{rem}
\source{intersection-complexes-unramified-l-factors}
\notready\label{rem:length}
\source{intersection-complexes-unramified-l-factors}
If $D$ is as above, and $\ch\lambda = \varrho_X(D)$, we can read off the length $\len (D)$ directly from $\ch\lambda$, if $X^\bullet$ happens to admit a $G$-eigen-volume form. Namely, under this assumption there is a character $\gamma\in \Lambda$ such that $\left< \gamma, \ch \nu_{D'}\right> = 1$ for every color $D'$, and therefore $\len (D) = \left< \gamma, \ch\lambda\right>$. Indeed, for such a color $D'$, there is a simple root $\alpha$ such that $D'P_\alpha$ contains the open Borel orbit, and then $D'P_\alpha \simeq \ch\nu_{D'}(\mathbb G_m)\backslash P_\alpha$. For such a homogeneous space to have a $P_\alpha$-eigen-volume form with eigencharacter $\mathfrak h$,  we must have that $\left< \mathfrak h+ 2\rho_{N_\alpha}, \ch\nu_{D'}\right> = 0$, where  $2\rho_{N_\alpha}$ is the sum of roots in the unipotent radical of $P_\alpha$. Equivalently, since $\left<\alpha,\ch\nu_{D'}\right>=1$, this reads $\left<\mathfrak h + 2\rho_G, \ch\nu_{D'}\right> = 1$. Therefore, 
\begin{equation}\label{e:length}
\source{intersection-complexes-unramified-l-factors}
 \len (D) = \left< \mathfrak h + 2\rho_G, \ch\lambda\right>,
\end{equation}
which we can unambiguously denote by $\len (\ch\lambda)$.
\end{rem}



\subsection{Proof of Theorem \ref{thm:comp}} 
\label{proof:comp}
\source{intersection-complexes-unramified-l-factors}

\subsubsection{Idea of the proof}


Let us give a set-theoretic idea for the proof of Theorem \ref{thm:comp}, as well as a guide to new notation we are about to introduce. The discussion here is not rigorous, using sets as avatars for geometric objects.

We have been using the notation $X^\bullet(F)_{G:\check\theta}$, or $(X^\bullet(F)/K)_{G:\check\theta}$ to denote the set of points in the $K$-orbit parametrized by $\ch\theta \in \mathcal V \cap \check\Lambda_X$, where $K = G(\mf o)$. The stratum  $\sM_X^{\check\theta}$ of the global model corresponds to $(X(\mathfrak o)/K)_{G:\check\theta}$, in the sense that it is determined by the condition that the ``$G$-valuation'' of maps in this stratum, at points where they fail to land in $X^\bullet$, is given by $\ch\theta$. 

Similarly, $T(\mathfrak o)N(F)$-orbits on $X^\circ(F)$ are parametrized by $\check \Lambda_X$, and let us denote by $X(F)_{B:\check\lambda}$ the subset of points which: (1) belong to $X^\circ (F)$, and (2) belong to the $T(\mathfrak o)N(F)$-orbit parametrized by $\check\lambda\in\check\Lambda_X$. As with the global model, the \emph{central fibers} of the stratum $\sY_X^{\check\lambda,\check\theta}$ of the Zastava space correspond to $(X(\mathfrak o)/B(\mathfrak o))_{\substack{G:\check\theta \\ B:\check\lambda}}$.


The following facts will be proven about the space $\sY_X^{\check\lambda,\check\theta}$:
\begin{itemize}
 \item It is nonempty only if $\check\lambda \succeq \check\theta$ (see Lemma \ref{lem:Ynonempty} and Corollary \ref{cor:Ytheta}). This is essentially a statement about the image of $X(\mathfrak o)$ in $X\sslash N(F)$. 
 \item It contains $\sY_X^{\check\lambda-\check\theta,0} \mathring\times \sY_X^{\check\theta,\check\theta}$ as an open dense (see Lemma \ref{lem:comp-cover}). This is a geometric statement, and it follows from the constructions that we describe below. 
\end{itemize}

Theorem \ref{thm:comp} says that the geometric analog of the action map 
\begin{equation}
 \label{eq:action-sets}
\source{intersection-complexes-unramified-l-factors}
\act: \overline{X^\bullet(\mathfrak o)} \times^{K}  \overline{K t^{\check\theta} K} \to \overline{X(\mathfrak o)}_{G:\check\theta}
\end{equation}
is birational and proper. (The closures are also understood here in the ``geometric'' topology, see Lemma \ref{lem:act-kpoints}.) 

To prove this, we fix $\check\theta \in \mathfrak c_X^-$, which will not appear in the notation, and work with $B(\mathfrak o)$-orbits, defining spaces whose central fibers satisfy the following set-theoretic analogies:
\begin{center}
\begin{tabular}{p{20mm}p{10mm}p{7cm}}
 $Z^{\check\lambda - \check\nu, \check\lambda}$ && $\overline{X^\bullet(\mathfrak o)}_{B:\check\lambda-\check\nu} \overset{B(\mathfrak o)}\times (T(\mathfrak o) t^{\check\nu} N(F) \cap \overline{K t^{\check\theta} K})/B(\mathfrak o)$ \\
 $Z^{\check\lambda - \check\nu, \check\lambda, \check\theta',\check\eta}$ && $\overline{X^\bullet(\mathfrak o)}_{\substack{G:\check\theta' \\ B:\check\lambda-\check\nu}} \overset{B(\mathfrak o)}\times (T(\mathfrak o) t^{\check\nu} N(F) \cap K t^{\check\eta} K)/B(\mathfrak o)$ \\ 
\end{tabular}
\end{center}

Here, $\ch\eta \ge \ch\theta$ (and antidominant), so that the Cartan double coset $K t^{\check\eta} K$ belongs (in the affine Grassmannian) to the closure of $K t^{\check\theta} K$. Both of the spaces above are subspaces of the preimage of $X(\mathfrak o)_{B:\check\lambda}$ in $\overline{X^\bullet(\mathfrak o)} \times^{K} \overline{K t^{\check\theta} K}$ --- in fact, substacks, in the appropriate setting. Indeed, stratifying $\overline{K t^{\check\theta} K}$ by the ``Mirkovi\'c--Vilonen (MV) cycles'' corresponding to its intersection with the horocycles $K t^{\check\nu} N(F)$, we obtain a stratification of this space by 
\[\left(\overline{X^\bullet(\mathfrak o)}  \times^K (K t^{\check\nu} N(F) \cap \overline{K t^{\check\theta} K})\right)/B(\mathfrak o) \cap \act^{-1} (X(\mathfrak o)/B(\mathfrak o))_{B:\check\lambda}\]
(where $\act$ denotes the action map). Note that $K t^{\check\nu} N(F) /B(\mathfrak o) = K\times^{B(\mathfrak o)} T(\mathfrak o) t^{\check\nu} N(F) /B(\mathfrak o)$, hence the above can also be written 
\[\left(\overline{X^\bullet(\mathfrak o)}  \times^{B(\mathfrak o)} (T(\mathfrak o) t^{\check\nu} N(F) \cap \overline{K t^{\check\theta} K})\right)/B(\mathfrak o) \cap \act^{-1} (X(\mathfrak o)/B(\mathfrak o))_{B:\check\lambda}.\]
An element of $\overline{X^\bullet(\mathfrak o)}$, multiplied by $t^{\check\nu} N(F)$, lands in $(X(\mathfrak o)/B(\mathfrak o))_{B:\check\lambda}$ if and only if it belongs to $\overline{X^\bullet(\mathfrak o)}_{B:\check\lambda-\check\nu}$, showing that the sets corresponding to $Z^{\check\lambda - \check\nu, \check\lambda}$ (and also to $Z^{\check\lambda - \check\nu, \check\lambda, \check\theta',\check\eta}$, after a further stratification by $G$-valuations) are indeed subspaces of $\overline{X^\bullet(\mathfrak o)} \times^{K} \overline{K t^{\check\theta} K}$. 


Here comes the important geometric input: the topology on the affine Grassmannian, resp.\ on the arc space of $X$ (or rather, its global models), implies that $Z^{\check\lambda - \check\theta, \check\lambda, 0,\check\theta}$ is open and dense. This is a combination of the statements that $X^\bullet(\mathfrak o)$ is (tautologically) open dense in $\overline{X^\bullet(\mathfrak o)}$ (allowing us to take $\check\theta'=0$), that the Cartan double coset $K t^{\check\theta} K$ is (again tautologically) open dense in its closure, and that the MV cycle $K t^{\ch\theta} K \cap K t^{\ch\theta} N(F)$ is open in it. 

The open MV stratum on the affine Grassmannian $K\backslash G(F)$ meets the coset $K\backslash K t^{\check\theta} K$ along the $N(\mathfrak o)$-orbit $K \backslash K t^{\ch\theta} N(\mf o)$. Therefore, the quotient
\[(T(\mathfrak o)t^{\check\theta} N(F) \cap K t^{\check\theta} K)/B(\mathfrak o)\]
is just a point ($\ch\theta$ is anti-dominant).
Thus, 
\[X^\bullet(\mathfrak o)_{B:\check\lambda-\ch\theta} \times^{B(\mathfrak o)} (T(\mathfrak o) t^{\check\theta} N(F) \cap K t^{\check\theta} K)/B(\mathfrak o) = (X^\bullet(\mathfrak o)/B(\mathfrak o))_{B:\check\lambda-\ch\theta},\]
which, on the other hand, can be identified as an open dense subset of $(X(\mathfrak o)/B(\mathfrak o))_{\substack{G:\ch\theta \\ B:\check\lambda}}$. Thus, in the geometric setting, the action map \eqref{eq:action-sets}, restricted to the subsets with $B$-valuation equal to some (any) $\check\lambda$, is indeed birational. On the other hand, the map $\sY^{\check\lambda}\to \sM_X$ is smooth for large $\check\lambda$ (Corollary \ref{cor:Mcover}), proving the birationality of (the geometric version of) \eqref{eq:action-sets}.








\subsubsection{}

We now turn to the actual proof of Theorem \ref{thm:comp}.

Since $\ol\Gr^{\ch\Theta}_{G,C^{\ch\Theta}}$ lives over $C^{\ch\Theta}$, 
we can factor $\act_\sM$ into 
\begin{equation} \label{e:actC}
\source{intersection-complexes-unramified-l-factors}
    \Scr M_X \ttimes \ol\Gr^{\ch\Theta}_{G,C^{\ch\Theta}} \overset{\act_{C^{\ch\Theta}}}\longrightarrow \Scr M_X \xt C^{\ch\Theta}
\overset{\pr_1}\longrightarrow \sM_X
\end{equation}
where $\act_{C^{\ch\Theta}}$ is the naturally induced map over $C^{\ch\Theta}$.


\begin{lem}
\source{intersection-complexes-unramified-l-factors}
\notready \label{lem:actimage} 
\source{intersection-complexes-unramified-l-factors}
Let $\ch\Theta \in \Sym^\infty(\mathfrak c_X^- \sm 0)$. 
\begin{enumerate}
\item
The preimage of $\Scr M^{\ch\Theta}_X$ in 
$\Scr M_{X} \ttimes \ol\Gr^{\ch\Theta}_{G,C^{\ch\Theta}}$ 
under $\act_\sM$ 
is contained in the open substack
$\Bun_H \ttimes \Gr^{\ch\Theta}_{G,\oo C^{\ch\Theta}}$. 

\item
The image $\act_\sM(\Bun_H \ttimes \ol\Gr^{\ch\Theta}_{G,\oo C^{\ch\Theta}}) \subset \Scr M_X$ contains $\Scr M_X^{\ch\Theta}$ as an open dense substack. 
\end{enumerate}
\end{lem}
\begin{proof}
Statement (i) follows from the description of
$\act_\sM$ on $k$-points, Lemma~\ref{lem:act-kpoints}, and the fact that no element of $\mathfrak c_X^- \sm 0$ is $ \preceq 0 $. 

To simplify notation, we give the proof of (ii) in the case when
 $\ch\Theta = [\ch\theta], \ch\theta\ne 0$ is a singleton partition. 
The general case is entirely analogous.

Now we are considering $\act_C : \sM_X \xt \ol\Gr^{\ch\theta}_{G,C} \to \sM_X \times C$. 
Recall that we have a 
locally closed embedding $\sM_X^{\ch\theta} \into \sM_X \times C$. 
First we show that 
$M:=\act_C(\Bun_H \ttimes \ol\Gr^{\ch\theta}_{G,C}) \subset \sM_X\times C$ contains 
$\Scr M^{\ch\theta}_X$.
Fix $v\in \abs C$. 
Take an arbitrary $H$-bundle $\sP_H \in \Bun_H(k)$, 
and choose a trivialization of $\sP_H|_{\Spec \mf o_v}$, which 
induces a trivialization $\tau_0 : (\sP_H \xt^H G)|_{\Spec \mf o_v} \cong \sP_G^0|_{\Spec \mf o_v}$. 
Then $(\sP_H,v,\tau_0)$ defines a point of $\wh\sM_X$. 
We also have the point $t^{\ch\theta} \in \Gr^{\ch\theta}_{G,v}$. 
The image of $(\sP_H,v,\tau_0, t^{\ch\theta})$ gives a point in 
$\sM_X \ttimes \Gr^{\ch\theta}_{G,v}$,
and by construction $\act_C$ will send this point to the stratum $\sM^{\ch\theta}_X$.

Thus we have shown that  
$M$ contains a point of $\sM^{\ch\theta}_X$ over every point $v\in \abs C$. 
It follows from Lemma~\ref{lem:act-equivariant} that 
$M \subset \Scr M_X \xt C$ 
is stable under
generic-Hecke modifications away from the marked point in $C$. 
Since these generic-Hecke modifications act transitively on 
the stratum $\Scr M^{\ch\theta}_X$ (Proposition~\ref{prop:genHecke-transitive}), we 
deduce that $M$ contains all of $\sM^{\ch\theta}_X$. 

The previous paragraph and Lemma~\ref{lem:act-kpoints} imply that 
\[ \sM_X^{\ch\theta} = M \sm \bigcup_{\ch\theta' \succ \ch\theta} \act_C(\sM_X \ttimes \ol\Gr^{\ch\theta'}_G), \] 
which shows that $\sM_X^{\ch\theta}$ is open in $M$. 


Since $\ol\Gr^{\ch\theta}_G$ is irreducible, 
for every connected component $U$ of $\Bun_H$, 
the image $M_U := \act_C(U \ttimes \ol\Gr^{\ch\theta}_{G,C})$ is 
irreducible. 
It is easy to see that its intersection with 
$\sM_X^{\ch\theta}$ is nonempty, hence dense. 
\end{proof}

As a corollary, we can now prove Theorem \ref{thm:comp}(i):

\begin{proof}[Proof of Theorem \ref{thm:comp}(i)]
Let $M$ denote an irreducible component of $\barM^0_X$. 
Then $M':=\act_\sM(M \ttimes \ol\Gr^{\ch\Theta}_{G,C})$
is an irreducible closed substack of $\Scr M_X$, 
and since $M$ contains a connected component of $\Bun_H$,
Lemma~\ref{lem:actimage}(ii) implies that 
 $M' \cap \Scr M^{\ch\Theta}_X$ is dense in $M'$.
It follows that $\Scr M^{\ch\Theta}_X$ is dense in $\act_\sM(\barM^0_X \ttimes \ol\Gr^{\ch\theta}_{G,C})$.
\end{proof}


\subsubsection{Base change to Zastava model}  
Fix $\ch\theta \in \mathfrak c_X^- \sm 0$ 
and a point $v \in \abs C$. 
Consider the restriction of $\act_\sM$ to 
$\sM_{X} \ttimes \ol\Gr^{\ch\theta}_{G,v} \to \sM_X$,
which is the fiber of $\act_C$ over 
$v \to C = C^{[\ch\theta]}$. 

Let us consider for $\ch\lambda\in \mathfrak c_X$ 
the fiber product diagram
\begin{equation}  \label{e:Zact-diagram}
\source{intersection-complexes-unramified-l-factors}
\begin{tikzcd}
Z_v^{?,\ch\lambda} \ar[r, "\act_{\Scr Y}"] \ar[d] & \Scr Y^{\ch\lambda} \ar[d] \\ 
\barM_X^0 \ttimes \ol\Gr^{\ch\theta}_{G,v} \ar[r, "\act_\sM"] & \Scr M_X 
\end{tikzcd} 
\end{equation}
An $S$-point of $Z_v^{?,\ch\lambda}$ consists of the data $(\sigma,\Scr P_G, \Scr P'_B, 
v, \tau)$ where 
\begin{itemize}
\item $(\sigma,\Scr P_G) \in \barM_X^0$,
\item $\Scr P'_B\in \Bun_B^{-\ch\lambda}$ is a $B$-structure on a $G$-bundle $\Scr P'_G := G\xt^B \Scr P'_B$, 
\item $\tau : \Scr P'_G|_{(C\sm v)\xt S} \cong \Scr P_G|_{(C\sm v)\xt S}$ is 
a modification inducing a point in $\ol\Gr^{\ch\theta}_G$ such that 
$\tau^{-1}\circ \sigma$ generically lands in $X^\circ \xt^B \Scr P'_B$. 
\end{itemize}

\noindent
Let $\widehat{\Scr Y} \to \Scr Y$
denote the Zariski locally trivial $\msf  L^+B$-torsor parametrizing 
$(\sigma,\Scr P_B)\in \Scr Y$ and a trivialization 
$\Scr P_B|_{\wh C'_v} \cong \Scr P_B^0|_{\wh C'_v}$, and let $\widehat{\Scr Y}_0 \to \Scr Y_0$ be the restriction to $\barM_X^0$. (We will use the index $0$ for the same purpose on the strata of $\Scr Y$.)


\begin{prop-const} \label{prop:Yact-strat}
\source{intersection-complexes-unramified-l-factors}
\ 
\begin{enumerate}
\item The fiber product $Z_v^{?,\ch\lambda}$ admits a stratification by 
\[
    Z_v^{\ch\lambda-\ch\nu,\ch\lambda} := 
    \widehat{\Scr Y}_0 \xt^{\msf L^+B} 
    (\msf L^+ T \cdot t^{\ch\nu}\cdot \msf L N \cap \ol{ \msf L^+G\cdot t^{\ch\theta} \cdot \msf L^+G}) / \msf L^+B, 
\]
where $\ch\nu$ ranges over the weights of the irreducible $\ch G$-module $V^{\ch\theta}$. 

\item
We have an isomorphism at the level of reduced schemes 
\[ Z_v^{\ch\lambda-\ch\nu,\ch\lambda} \cong 
\Scr Y_0^{\ch\lambda-\ch\nu} \ttimes 
(\msf S^{\ch\nu} \cap \ol\Gr^{\ch\theta}_G). \] 
where $\Scr Y_0^{\ch\lambda-\ch\nu} \ttimes -$ 
means 
$\widehat{\Scr Y}_0
\xt^{\msf L^+B} -$. 

\item
The open stratum corresponds to $\ch\nu=\ch\theta$. More precisely, we have
\[ Z_v^{\ch\lambda-\ch\theta,\ch\lambda} \cong \Scr Y_0^{\ch\lambda-\ch\theta} \xt v, \] 
where $v$ corresponds to the embedding $\{ t^{\ch\theta} \} 
\into \Gr^{\ch\theta}_{G,v}$.
\end{enumerate}
\end{prop-const}



\begin{proof}
Recall that $\ol{\msf L^+G\cdot t^{\ch\theta}\cdot \msf L^+G}$ has a stratification 
by intersecting with $\msf L^+G\cdot t^{\ch\nu}\cdot \msf L N$. 
If we take the quotient on the left by $\msf L^+ G$, then 
under the identification 
$\msf L^+G \bs \msf L G \cong \Gr_G : g \mapsto g^{-1}$, we have
described above the stratification of $\ol\Gr^{-\ch\theta}_G$ 
by MV cycles $\ol\Gr^{-\ch\theta}_G \cap \msf S^{-\ch\nu}$, where
$-\ch\theta$ is dominant. 
It is known by \cite[Theorem 3.2]{MV} that 
$\ol\Gr^{-\ch\theta}_G \cap \msf S^{-\ch\nu}$
is non-empty precisely when $-\ch\nu$ is a weight of 
$V^{-\ch\theta}$, and the open stratum corresponds to 
when $-\ch\nu$ equals $-\ch\theta$.

Now let $(\sigma,\Scr P_G, \Scr P'_B,\tau)$ be an $S$-point of 
$Z_v^{?,\ch\lambda}$. The restriction of $(\Scr P_G,\Scr P'_B,\tau)$
to $\wh C'_v$ gives a point in $\msf L^+G \bs \ol{\msf L^+G \cdot t^{\ch\theta} \cdot \msf L^+ G} / \msf L^+ B$. 
Therefore, by the previous paragraph, we can stratify $Z_v^{?,\ch\lambda}$
by the preimages of 
\[  \msf L^+G \bs (\ol{\msf L^+G\cdot t^{\ch\theta}\cdot \msf L^+G} \cap \msf L^+G\cdot t^{\ch\nu}\cdot \msf L N) / \msf L^+B . \] 
Suppose our $S$-point lies in such a stratum corresponding to $\ch\nu$. 
In particular, $\tau$ 
corresponds to a point in $\msf L^+B \bs (\Gr^{\ch\nu}_B)_\red$,
which means that there exists a $B$-structure $\Scr P_B|_{\wh C'_v}$
on $\Scr P_G|_{\wh C'_v}$ such that $\tau$ gives an isomorphism
of generic $B$-bundles
$\tau : \Scr P'_B|_{\wh C^\circ_v} \cong \Scr P_B|_{\wh C^\circ_v}$.  
By Beauville--Laszlo's theorem, the datum $(\Scr P'_B|_{C\sm v}, \Scr P_B|_{\wh C'_v},\tau)$ descends to a $B$-structure $\Scr P_B$ on $\Scr P_G$
such that 
$\Scr P'_B|_{C\sm v} \cong \Scr P_B|_{C\sm v}$. 
Then $(\sigma, \Scr P_B) \in \sY_0^{\ch\lambda-\ch\nu}(S)$ and 
$(\sigma,\Scr P_B, \Scr P'_B, v, \tau) \in Z_v^{\ch\lambda-\ch\nu,\ch\lambda}(S)$. 
The procedure above can be reversed to see that $Z_v^{\ch\lambda-\ch\nu,\ch\lambda}$ is equal to the entire stratum.
This shows (i). 

\smallskip

Observe that $\msf L^+ T \cdot t^{\ch\nu} \cdot \msf L N = \msf L N \cdot t^{\ch\nu} \cdot \msf L^+B$ and 
there is a map from 
$(\msf L N \cdot t^{\ch\nu} \cdot \msf L^+B \cap \ol{\msf L^+G \cdot t^{\ch\theta} \cdot \msf L^+G})/\msf L^+B$ to $\msf S^{\ch\nu}\cap \ol\Gr^{\ch\theta}_G$ which is an isomorphism at the level of reduced schemes. 
Now (ii) follows from (i).

The open stratum $Z_v^{\ch\lambda-\ch\theta,\ch\lambda}$ 
corresponds to the open stratum $\ol\Gr^{-\ch\theta}_G \cap \msf S^{-\ch\theta}$, 
when $\ch\nu=\ch\theta$. 
By \cite[(3.6)]{MV}, we have $\msf S^{\ch\theta} \cap \ol\Gr^{\ch\theta}_G = \{t^{\ch\theta}\}$
so (iii) is a special case of (ii).
\end{proof}

Observe that $\barM_X^0 \ttimes \ol\Gr^{\ch\theta}_{G,v}$ has 
a stratification by $\sM_{X}^{\ch\Theta'} \ttimes \Gr^{\ch\eta}_{G,v}$ for those $\ch\Theta' \in \Sym^\infty(\mathfrak c_X^- \sm 0)$ such that the stratum  $\sM_{X}^{\ch\Theta'}$ belongs to $\barM_X^0$ --- still to be determined, see Lemma \ref{lem:M0closure} --- 
and those $\ch\eta \in \ch\Lambda^-_G$ such that $\ch\eta \ge \ch\theta$ (equivalently, $\ch\eta$ is a weight of $V^{\ch\theta}$). 
Let $Z^{?,\ch\lambda,\ch\Theta',\ch\eta}_v$ denote the
preimage of the corresponding stratum in $Z^{?,\ch\lambda}_v$, 
so we have a Cartesian square 
\begin{equation} \label{e:Z0act-diagram}
\source{intersection-complexes-unramified-l-factors}
\begin{tikzcd} 
Z_v^{?,\ch\lambda,\ch\Theta',\ch\eta} \ar[r, "\act_\sY"] \ar[d] & \Scr Y^{\ch\lambda}_X \ar[d] \\ 
\sM^{\ch\Theta'}_{X} \ttimes \Gr^{\ch\eta}_{G,v} \ar[r, "\act_\sM"] & \Scr M_X 
\end{tikzcd}
\end{equation}
This diagram now has no dependence on $\ch\theta$ and is defined
entirely with respect to $\ch\eta$. 
Proposition~\ref{prop:Yact-strat} implies that there is a
stratification
\[ Z_v^{?,\ch\lambda,\ch\Theta', \ch\eta} = \bigcup_{\ch\nu} Z_v^{\ch\lambda-\ch\nu,\ch\lambda,\ch\Theta', \ch\eta} \]
where $\ch\nu$ runs through the weights of $V^{\ch\eta}$, and  
$Z_v^{\ch\lambda-\ch\nu,\ch\lambda,\ch\Theta',\ch\eta}$ 
admits a map to $\Scr Y^{\ch\lambda-\ch\nu,\ch\Theta'}$. 
The open stratum is
\[ Z_v^{\ch\lambda-\ch\eta,\ch\lambda,\ch\Theta',\ch\eta} \cong 
\Scr Y^{\ch\lambda-\ch\eta,\ch\Theta'} \times \{ t^{\ch\eta}\}. \]

We will make special use of the case $\ch\Theta'=0$ and $\ch\eta = \ch\theta$ 
because 
Lemma~\ref{lem:actimage} implies that $\Scr Y^{\ch\lambda,\ch\theta}$
is contained in the images of $Z_v^{?,\ch\lambda,0,\ch\theta} \to \Scr Y^{\ch\lambda}$ 
over all $v \in \abs C$. 
Note that $Z_v^{\ch\lambda-\ch\theta,\ch\lambda,0,\ch\theta} = \sY^{\ch\lambda-\ch\theta,0}$.

\smallskip 

Recall that $\mathfrak c_X^{\Cal D}$ denotes the monoid
generated by $\ch\nu_D$ for $D\in \Cal D$. 
\begin{cor}
\source{intersection-complexes-unramified-l-factors}
\notready \label{cor:Ytheta}
\source{intersection-complexes-unramified-l-factors}
Let $\ch\theta \in \mathfrak c_X^- \sm 0$. 
\begin{enumerate}
\item
The scheme $\Scr Y^{\ch\lambda,\ch\theta}$ is nonempty only if 
$\ch\lambda \succeq \ch\theta$. 

\item
We have an isomorphism $(\Scr Y^{\ch\theta,\ch\theta})_\red \cong C$ over
the diagonal $C \into \Scr A^{\ch\theta}$.

\item
The single point in the central fiber 
$\msf Y^{\ch\theta,\ch\theta}$ 
corresponds to $t^{\ch\theta} \in \msf S^{\ch\theta}\cap \Gr^{\ch\theta}_G$. 
\end{enumerate}
\end{cor}
\begin{proof}
As we remarked above, every point of $\sY^{\ch\lambda,\ch\theta}$ is contained
in the image of $Z_v^{?,\ch\lambda,0,\ch\theta}$ for some $v\in \abs C$. 
By the description of the stratification
of $Z_v^{?,\ch\lambda,0,\ch\theta}$, the latter is nonempty only
if $\Scr Y^{\ch\lambda-\ch\nu,0}$ is nonempty for some $\ch\nu$. 
This implies that $\ch\lambda \succeq \ch\nu$ by 
Lemma~\ref{lem:Ynonempty}. Since 
$\ch\nu \ge \ch\theta$, we deduce (i). 

If $\ch\lambda=\ch\theta$, then $\Scr Y^{\ch\theta-\ch\nu,0}$ 
is only nonempty when $\ch\nu=\ch\theta$, in which case $\Scr Y^0 = \pt$ and $Z_v^{0,\ch\theta,0,\ch\theta} = \{t^{\ch\theta}\}$. 
By moving the point $v$ around, we get 
a surjection $C \to (\Scr Y^{\ch\theta,\ch\theta})_\red$ and
it must be an isomorphism (on reduced schemes) since the composition 
with $\pi :\Scr Y^{\ch\theta,\ch\theta} \to \Scr A^{\ch\theta}$ gives 
the diagonal embedding.
\end{proof}

Let $\Scr Y^{\ch\lambda-\ch\theta,0} \oo\xt v$ denote the disjoint locus 
$\Scr Y^{\ch\lambda-\ch\theta,0} \oo\xt \msf Y^{\ch\theta,\ch\theta}_v$. 
Observe that the composition
\[ \sY^{\ch\lambda-\ch\theta} \oo\xt v \into 
\sY^{\ch\lambda-\ch\theta} \xt v = 
Z^{\ch\lambda-\ch\theta,\ch\lambda}_v \overset{\act_\sY}\longrightarrow
\sY^{\ch\lambda} \]
coincides with the composition 
$\sY^{\ch\lambda-\ch\theta} \oo\xt v \into 
\sY^{\ch\lambda-\ch\theta} \oo\xt \sY^{\ch\theta,\ch\theta} \to 
\sY^{\ch\lambda}$ 
where the second map is given by the graded factorization property.


\begin{lem}
\source{intersection-complexes-unramified-l-factors}
\notready \label{lem:comp-cover}
\source{intersection-complexes-unramified-l-factors}
Let $\ch\theta\in \mathfrak c_X^- \sm 0$. 
Then:
\begin{enumerate}
\item The open embedding 
$ \sY^{\ch\lambda-\ch\theta,0} \oo\xt \sY^{\ch\theta,\ch\theta} \into
\sY^{\ch\lambda,\ch\theta}$
given by the graded factorization property is dense. 
\item For $\ch\lambda$ large enough, 
the open stratum $\sY^{\ch\lambda-\ch\theta,0} \xt v
\into Z^{?,\ch\lambda,0,\ch\theta}_v$ is dense. 
\end{enumerate}
\end{lem}
\begin{proof}
First observe that by Lemma~\ref{lem:localglobalyoga} we may
always assume that $\ch\lambda$ is large enough so that
the conditions of Corollary~\ref{cor:Mcover} hold: namely, 
$\sY^{\ch\lambda}\to \sM_X$ is smooth with connected fibers. 

By definition, $\sY^{\ch\lambda,\ch\theta}$ is the preimage 
of $\sM^{\ch\theta}_X$. Now let $V$ be a connected component of 
$\sM^{\ch\theta}_X$ that intersects the image of $\sY^{\ch\lambda}$. 
Then 
$\sY^{\ch\lambda,\ch\theta}_V := \sY^{\ch\lambda} \xt_{\sM_X} V$ is
a nonempty connected component of the smooth scheme
$\sY^{\ch\lambda,\ch\theta}$.
Lemma~\ref{lem:actimage} implies that $V$ lies in  
$\act_\sM(U \ttimes \Gr^{\ch\theta}_{G,C})$, where $U$ is a connected
component of $\Bun_H$. 
Consider the fiber product diagram 
\begin{equation}  \label{e:Z0f-diagram}
\source{intersection-complexes-unramified-l-factors}
\begin{tikzcd} 
Z^{?,\ch\lambda,0,\ch\theta} \ar[r]\ar[d] & \sY^{\ch\lambda} \ar[d] \\ 
\Bun_H \ttimes \Gr^{\ch\theta}_{G,C} \ar[r] & \sM_X 
\end{tikzcd}
\end{equation}
which is the analog of \eqref{e:Z0act-diagram} where we allow $v$
to vary. 
The fibers of $\sY^{\ch\lambda}\to \sM_X$ are irreducible, so
$Z^{?,\ch\lambda,0,\ch\theta}_U := Z^{?,\ch\lambda,0,\ch\theta} \xt_{\Bun_H} U$ is irreducible and the image of 
\[ Z^{?,\ch\lambda,0,\ch\theta}_U \to \sY^{\ch\lambda} \]
contains $\sY^{\ch\lambda,\ch\theta}_V$ as a dense open. 
Proposition~\ref{prop:Yact-strat} (twisted by $C$) implies that 
$Z^{?,\ch\lambda,0,\ch\theta}_U$ has a stratification 
$\bigcup_{\ch\nu} Z^{\ch\lambda-\ch\nu,\ch\lambda,0,\ch\theta}_U$ 
over the weights $\ch\nu$ of $V^{\ch\theta}$, and 
$Z_U^{\ch\lambda-\ch\nu,\ch\lambda,0,\ch\theta}$ maps to 
$\sY^{\ch\lambda-\ch\nu,0}\xt_{\Bun_H} U$. 
In particular, $\sY^{\ch\lambda-\ch\nu,0}_U := \sY^{\ch\lambda-\ch\nu,0}\xt_{\Bun_H} U$ is nonempty
for some $\ch\nu$. Since $\ch\nu \ge \ch\theta$, 
Corollary~\ref{cor:Mcover}(ii) implies that $\sY^{\ch\lambda-\ch\theta,0}_U
$ is also nonempty.
Therefore, the open stratum 
$Z_U^{\ch\lambda-\ch\theta,\ch\lambda,0,\ch\theta} \cong 
\sY^{\ch\lambda-\ch\theta,0}_U \times C$ is nonempty. 
By Corollary~\ref{cor:Ytheta}, we can identify 
$C \cong (\sY^{\ch\theta,\ch\theta})_\red$, and  
the map 
\[ \sY^{\ch\lambda-\ch\theta,0} \oo\xt C \into 
Z^{?,\ch\lambda,0,\ch\theta} \to \sY^{\ch\lambda} \]
coincides with the open embedding 
$\sY^{\ch\lambda-\ch\theta,0} \oo\xt \sY^{\ch\theta,\ch\theta} \into \sY^{\ch\lambda,\ch\theta}$ 
given by graded factorization. 
Then $\sY^{\ch\lambda-\ch\theta}_U \oo\xt C$ is a nonempty 
open subscheme of the irreducible component $\sY^{\ch\lambda,\ch\theta}_V$,
so it must be dense. 

To show (ii), note that \emph{a priori}  
the preimage of 
$\sY^{\ch\lambda,\ch\theta}_V$ in $Z^{?,\ch\lambda,0,\ch\theta}$
is contained in a finite union of 
$Z^{?,\ch\lambda,0,\ch\theta}_U$
for connected components $U$ of $\Bun_H$.
However we saw above that the open stratum $\sY^{\ch\lambda-\ch\theta,0}_U \times C$ is nonempty, and therefore dense in $Z^{?,\ch\lambda,0,\ch\theta}_U$, 
for each such $U$. 
\end{proof}

In the proof above, we have essentially shown 
Theorem \ref{thm:comp}(ii) along the way:

\begin{proof}[Proof of Theorem \ref{thm:comp}(ii)] 
To simplify notation, we only show the case $\ch\Theta = [\ch\theta]$. 
The multi-point version is proved in exactly the same way. 
We continue to use the notation from the previous proof
of Lemma~\ref{lem:comp-cover}. 
By Corollary~\ref{cor:Mcover} and the base change diagram \eqref{e:Z0f-diagram}, 
it is enough to show, by restricting to open dense subsets, that 
\[ Z^{?,\ch\lambda,0,\ch\theta} \to \sY^{\ch\lambda} \xt_{\sM_X} \barM_X^{\ch\theta} \] 
is birational for $\ch\lambda$ large enough. 
It follows from Lemma~\ref{lem:comp-cover} that 
$\sY^{\ch\lambda-\ch\theta,0} \oo\xt C$ is a dense open 
in both the target and source. 
\end{proof}


\begin{proof}[Proof of Theorem \ref{thm:comp}(iii)]
Again to simplify notation we only show the case $\ch\Theta = [\ch\theta]$.
By Zariski's Main Theorem, it suffices to show that 
the restriction of $\act_\sM$ to 
\begin{equation} \label{e:act-bijection}
\source{intersection-complexes-unramified-l-factors}
\act_\sM^{-1}(\sM^{\ch\theta}_X) \cap (\Bun_H \ttimes \Gr_{G,C}^{\ch\theta}) \to \sM_X^{\ch\theta}
\end{equation}
is a bijection on $k$-points. 
By Lemma~\ref{lem:act-equivariant}, this map is equivariant with respect to
generic-Hecke modifications away from the marked point in $C$. 
Fix a marked point $v\in C$ and let $\sM^{\ch\theta}_{X,v}$ denote 
the preimage of $v$ under the projection $\sM^{\ch\theta}_X \to C$.
By Proposition~\ref{prop:genHecke-transitive}, all the $k$-points 
of $\sM^{\ch\theta}_{X,v}$ are equivalent under 
the equivalence relation generated by generic-Hecke correspondences.  
Thus it suffices to determine the fiber of \eqref{e:act-bijection}
at the single point $\{t_v^{\ch\theta}\} = \msf Y^{\ch\theta,\ch\theta}$.  
By Proposition \ref{prop:Yact-strat}, 
the fiber of $t_v^{\ch\theta}$ has a stratification by 
$\msf Y^{\ch\theta-\ch\nu,0} \ttimes (\msf S^{\ch\nu} \cap \ol\Gr^{\ch\theta}_G)$
for $\ch\nu$ a weight of $V^{\ch\theta}$. 
On the other hand, Corollary \ref{cor:Ytheta} implies that 
$\msf Y^{\ch\theta-\ch\nu,0}$ is nonempty only if $\ch\nu=\ch\theta$. 
In this case $\msf Y^{0,0} \times (\msf S^{\ch\theta}\cap \ol\Gr^{\ch\theta}_G)$
is a point. 
\end{proof}

\subsubsection{}
Recall that for an arbitrary algebraic group $H$, the algebraic fundamental
group $\pi_1(H)$ is defined as the quotient of the coweight lattice by the coroot lattice
of the reductive group $H/H_{\mathrm u}$, where $H_{\mathrm u}$ is the unipotent
radical of $H$.
For $\ch\mu\in \pi_1(H)$, let $\Bun_H^{\ch\mu}$ denote the
corresponding connected component of $\Bun_H$ and let $^{\ch\mu}\barM_X^0$ 
be its closure in $\sM_X$.

\begin{cor}
\source{intersection-complexes-unramified-l-factors}
\notready \label{cor:strata-comp} 
\source{intersection-complexes-unramified-l-factors}
Let $\ch\Theta \in \Sym^\infty( \mathfrak c_X^- \sm 0)$. 
Then the set of irreducible components of $\barM^{\ch\Theta}_X$ is in bijection 
with $\pi_1(H)$, where $\ch\mu\in \pi_1(H)$ corresponds to 
\[ ^{\ch\mu}\barM_X^{\ch\Theta} := 
\act_\sM(^{\ch\mu}\barM^0_X \ttimes \ol\Gr^{\ch\Theta}_{G,C^{\ch\Theta}}). 
\] 
\end{cor}

\begin{proof}
 It follows from Theorem \ref{thm:comp} that the irreducible
components of $\barM^{\ch\Theta}_X$ are in bijection with the
irreducible components of $\Bun_H \ttimes \ol\Gr^{\ch\Theta}_{G,C^{\ch\Theta}}$. 
Since $\ol\Gr^{\ch\Theta}_{G,C^{\ch\Theta}}$ is irreducible, the
latter are in bijection with $\pi_0(\Bun_H) = \pi_1(H)$.
\end{proof}

We will let $^{\ch\mu}\sM_X^{\ch\Theta} := {^{\ch\mu}}\barM_X^{\ch\Theta} \cap \sM_X^{\ch\Theta}$
denote the corresponding connected component of $\sM_X^{\ch\Theta}$ (which is smooth). 

\subsubsection{Mirkovi\'c--Vilonen cycles} 

We finish this subsection by proving a result that will be used in the following sections. The goal here is to show that the Mirkovi\'c--Vilonen cycles in the $\ch\theta$-stratum of the affine Grassmannian map, generically, to the $\ch\theta$-stratum of the global model of $X$, under the action map.

Fix $v\in \abs C$ and $\ch\theta \in \mf c_X^- \sm 0$. 
Consider the restriction of $\act_{\sM,v}$ to 
\begin{equation} \label{e:heckemapGr}
\source{intersection-complexes-unramified-l-factors}
    \pt \xt \ol\Gr^{\ch\theta}_{G,v} \to \Bun_H \ttimes \ol\Gr^{\ch\theta}_{G,v} 
\to \barM_X^{\ch\theta}, 
\end{equation}
where $\pt \to \Bun_H$ corresponds to the trivial $H$-bundle. 
Note that the map \eqref{e:heckemapGr} above can be 
extended to a map 
\begin{equation}\label{e:actHG}
\source{intersection-complexes-unramified-l-factors}
    \act_{v} : \Gr_G \xt_{\msf L X/\msf L^+G} (\msf L^+X/\msf L^+G) \to \sM_X
\end{equation}
using Beauville--Laszlo's theorem: a point of the left hand side 
consists of a $G$-bundle $\Scr P_G$ and a trivialization 
$\tau : \sP_G|_{C\sm v} \cong \sP_G^0|_{C\sm v}$ such that 
$\tau^{-1} \circ x_0 : C\sm v \to X \xt^G \sP_G$ is regular when localized at $\mf o_v$. 
Here $x_0 : C \to X \xt^G \Scr P_G^0 = X \xt C$ denotes the section corresponding to
the base point $x_0\in X^\circ$.  
Thus, $\act_{v}(\sP_G,\tau) := (\sP_G, \tau^{-1}\circ x_0) \in \sM_X^{\ch\theta}$
is well-defined. 

Define $\overset\circ\Gr{}^{\ch\theta}_G \subset \Gr^{\ch\theta}_G$ to be the 
open subscheme equal to the preimage 
of the stratum $\sM_X^{\ch\theta}$ under \eqref{e:heckemapGr}. 
We can also identify
$\displaystyle \overset\circ\Gr{}_G^{\ch\theta} = \Gr^{\ch\theta}_G \xt_{\msf L X/\msf L^+G} (\msf L^{\ch\theta}X / \msf L^+G)$.


Taking central fibers with respect to $v$, 
the restriction of $\act_v$ to a semi-infinite orbit
factors through 
\begin{equation}
 \label{e:actvY}
\source{intersection-complexes-unramified-l-factors}
 \msf S^{\ch\lambda} \cap \ol\Gr^{\ch\theta}_G \into \msf Y^{\ch\lambda}\xt_{\sM_X} \barM_X^{\ch\theta}, 
\end{equation}
where $\ch\lambda$ is a weight of $V^{\ch\theta}$
and we consider $(\msf Y^{\ch\lambda})_\red$ as a subscheme of $\msf S^{\ch\lambda}$ via
Lemma~\ref{lem:Y=ScapGr}.
Note that $\msf S^{\ch\lambda}\cap\ol\Gr_G^{\ch\theta}$ is isomorphic to 
the \emph{closed} stratum $Z_v^{0,\ch\lambda}$ from Proposition~\ref{prop:Yact-strat}.
Under this identification \eqref{e:actvY} coincides with the restriction of the map $\act_\sY:Z_v^{?,\ch\lambda}\to \sY^{\ch\lambda}$ from \eqref{e:Zact-diagram}.

Observe that $\msf S^{\ch\theta} \cap \ol\Gr^{\ch\theta}_G = \{ t^{\ch\theta} \}$ 
is contained in $\overset\circ{\Gr}{}^{\ch\theta}_G$. 
We will use this to deduce: 

\begin{lem}
\source{intersection-complexes-unramified-l-factors}
\notready \label{lem:MVcentralfiber}
\source{intersection-complexes-unramified-l-factors}
Let $\ch\lambda,\ch\theta$ as above. Then 
$\msf S^{\ch\lambda} \cap \overset\circ\Gr{}^{\ch\theta}_G$
intersects every irreducible component 
of $\msf S^{\ch\lambda} \cap \ol\Gr^{\ch\theta}_G$. 
\end{lem}
\begin{proof}
Let $Z$ denote an irreducible component of $\msf S^{\ch\lambda}\cap \ol\Gr^{\ch\theta}_G$, which must be 
of dimension $d=\brac{\rho_G, \ch\lambda-\ch\theta}$.
Let $\ol Z$ denote its closure in $\olsf S^{\ch\lambda} \cap \ol\Gr^{\ch\theta}_G$. 
Then the proof of \cite[Theorem 3.2]{MV}, which we briefly recall in the next paragraph,
shows that $\ol Z$ contains $t^{\ch\theta} \in \overset\circ\Gr{}^{\ch\theta}_G$.
Thus, $\ol Z \cap \overset\circ\Gr{}^{\ch\theta}_G$ is open and nonempty, hence dense 
in $\ol Z$. 

Since the boundary of $\msf S^{\ch\lambda}$ in $\olsf S^{\ch\lambda}$ 
is a hyperplane section (Proposition~\ref{prop:Shyperplane}), 
$\ol Z\sm Z$ contains an irreducible component of dimension $d-1$
inside $\msf S^{\ch\lambda_1}\cap \ol\Gr^{\ch\theta}_G$ for $\ch\lambda_1 < \ch\lambda$. 
In this way we produce a sequence $\ch\lambda=\ch\lambda_0, \dotsc, \ch\lambda_d$
and irreducible components of $\msf S^{\ch\lambda_i} \cap \ol Z$ 
of dimension $d-i$. 
The only weight $\ch\lambda_d$ of $V^{\ch\theta}$ such that 
$\brac{\rho_G,\ch\lambda-\ch\lambda_d} \ge d$ 
is $\ch\lambda_d=\ch\theta$, so $\ol Z$ must contain $\msf S^{\ch\theta}\cap \ol\Gr^{\ch\theta}_G
= \{ t^{\ch\theta} \}$. 
\end{proof}



\subsection{Closure relations and components in the global model}
\label{sect:closures}
\source{intersection-complexes-unramified-l-factors}

Let $\ch\Theta, \ch\Theta' \in \Sym^\infty(\mathfrak c_X^- \sm 0)$. Consider $\ch\Theta'-\ch\Theta$ as a formal sum 
in $\bigoplus_{\ch\theta\in \mathfrak c_X^-\sm 0} \mbb Z [\ch\theta]$. We say that 
\[ \ch\Theta' \succeq \ch\Theta \] 
if $\ch\Theta'-\ch\Theta$ can be written as a sum of formal differences 
$[\ch\theta']-[\ch\theta]$ where $\ch\theta,\ch\theta'\in \mathfrak c_X^-$ and $\ch\theta' \succeq \ch\theta$. 
Note that we allow $\ch\theta=0$ (in which case $[0]=0$), 
and for general $\ch\theta,\ch\theta'$ as
above, it is \emph{not} necessarily the case that $\ch\theta-\ch\theta' \in \Cal V$. 

\begin{prop}
\source{intersection-complexes-unramified-l-factors}
\notready \label{prop:closure-rel}
\source{intersection-complexes-unramified-l-factors}
Let $\ch\Theta,\ch\Theta' \in \Sym^\infty(\mathfrak c_X^- \sm 0)$. 
We have that $\sM^{\ch\Theta'}_X$ lies in the closure of $\sM^{\ch\Theta}_X$
if and only if there exists $\ch\Theta'' \in \Sym^\infty(\mathfrak c_X^-\sm 0)$ such that $\ch\Theta$ refines $\ch\Theta''$ and $\ch\Theta' \succeq \ch\Theta''$. 
\end{prop}

\begin{rem}
\source{intersection-complexes-unramified-l-factors}
\notready
We warn that our notation in Corollary~\ref{cor:strata-comp} is unfortunately 
\emph{not} compatible
with the closure relations in the following sense: 
it is possible that $^{\ch\mu} \barM^{\ch\Theta}_X \cap \sM^{\ch\Theta'}_X 
\ne {^{\ch\mu}\sM^{\ch\Theta'}_X}$ for $\ch\Theta' \succ \ch\Theta$ 
and $\ch\mu \in \pi_1(H)$.
\end{rem}



The reader may want to skip the proof of this proposition at first reading, and focus on the corollaries that follow. Note that the proof will use Zastava models, despite the fact that the statement is about the global model. The proof of the proposition starts with the following special case:


\begin{lem}
\source{intersection-complexes-unramified-l-factors}
\notready \label{lem:M0closure}
\source{intersection-complexes-unramified-l-factors}
The closure of $\sM_X^0$ in $\sM_X$ intersects 
$\sM_X^{\ch\Theta}$ only if $\ch\Theta \succeq 0$. 
\end{lem}

\begin{proof}
Let $\ch\Theta$ correspond to a stratum such that 
$\sM_X^{\ch\Theta}$ intersects $\barM_X^0$. Then by
Corollary~\ref{cor:Mcover} there exists $\ch\lambda$ 
such that $\sY^{\ch\lambda,\ch\Theta}$ intersects the closure
of $\sY^{\ch\lambda,0}$ in $\sY^{\ch\lambda}$. 
Now consider the torsor $H'\bs G' \to X^\bullet$ from 
Lemma~\ref{lem:Xcover}. 
Let $X' = \Spec k[H'\bs G']$ considered as an affine $G'$-variety.
The corresponding compactified Zastava model 
$\barY^{D}_{X'} \to \sA^{D}_{X'}$ is indexed
by $D\in \mf c_{X'}=\mbb N^{\Cal D}$. 
Since $X'^\bullet/B' \cong X^\bullet/B$ as stacks, 
$\sY^{\ch\lambda,0}_X$ is a disjoint union of $\sY^D_{X^\bullet}=\sY^{D,0}_{X'}$
ranging over all $D\in \mbb N^{\Cal D}$ such that $\varrho_X(D)=\ch\lambda$
(see \S\ref{sect:Y0cc}). 
Choose $D\in \mbb N^{\Cal D}$ such that the closure of 
$\sY^D_{X^\bullet}$ in $\sY^{\ch\lambda}_X$ intersects 
the stratum $\sY^{\ch\lambda,\ch\Theta}_X$. 
The map 
\begin{equation} \label{e:mapbarYupdown}
\source{intersection-complexes-unramified-l-factors}
 \barY^{D}_{X'} \to \barY^{\ch\lambda}_X  
\end{equation}
is proper because $\barY^{D}_{X'}, \barY^{\ch\lambda}_X$
are proper over $\sA^{D}_{X'},\sA^{\ch\lambda}_X$, respectively,
and the natural map $\sA^{D}_{X'}\to \sA^{\ch\lambda}_X$ 
is proper. Therefore, the closure of $\sY^{D}_{X^\bullet}$
in $\sY^{\ch\lambda}_X$ is contained in the image
of \eqref{e:mapbarYupdown}.
Note that the stratification of $\sM_{X'}$ is indexed by 
$\ch\Theta' \in \Sym^\infty(\mathfrak c_{X'}^-\sm 0)$. 
In particular, $\ch\Theta' \succeq 0$ since $\mf c_{X'}=\mbb N^{\Cal D}$. 
Therefore, $\barY^{D}_{X'}$ is a union of 
$\barY^{D}_{X'} \xt_{\sM_{X'}} \sM_{X'}^{\ch\Theta'}$
for $\ch\Theta' \succeq 0$, which implies its image in 
$\barY^{\ch\lambda}_X$ is contained in the union 
of $\barY^{\ch\lambda}_X \xt_{\sM_X} \sM_X^{\ch\Theta}$
for $\ch\Theta \succeq 0$.
\end{proof}



\begin{proof}[{Proof of Proposition~\ref{prop:closure-rel}}] 
By Theorem \ref{thm:comp}(i), the closure of $\sM^{\ch\Theta}_X$ is equal to 
the image of $\barM^0_X \ttimes \ol\Gr^{\ch\Theta}_{G,C^{\ch\Theta}}$,
so we will consider the latter. 
Note that $C^{\ch\Theta}$ is stratified by disjoint loci $\oo C^{\ch\Theta''}$ 
for all partitions $\ch\Theta''$ such that $\ch\Theta$ refines 
$\ch\Theta''$. By the description of the fibers of 
$\Gr^{\ch\Theta}_{G, C^{\ch\Theta}} \to C^{\ch\Theta}$ in 
\eqref{e:GrTheta-fiber}, we have an identification 
\[
\ol\Gr^{\ch\Theta}_{G,C^{\ch\Theta}} \xt_{C^{\ch\Theta}} \oo C^{\ch\Theta''} 
= \ol\Gr^{\ch\Theta''}_{G, \oo C^{\ch\Theta''} } 
\]
at the level of reduced schemes. Therefore, replacing $\ch\Theta$ by $\ch\Theta''$, 
it suffices to show that 
the image of $\barM_X^0 \ttimes \ol\Gr^{\ch\Theta}_{G, \oo C^{\ch\Theta}}$
contains $\sM_X^{\ch\Theta'}$ if and only if $\ch\Theta' \succeq \ch\Theta$.

The ``only if'' direction follows from the description of
$\act_\sM$ on $k$-points and Lemmas~\ref{lem:act-kpoints} 
and \ref{lem:M0closure}. 

\medskip

We will show the ``if'' direction only 
in the case when $\ch\Theta=[\ch\theta],\, \ch\Theta'=[\ch\theta']$ 
are singleton (so $\ch\theta'\succeq \ch\theta$) to lessen notation 
(allowing $\ch\theta=0$), but the multi-point version is proved in 
exactly the same way. 
Fix $v\in \abs C$. We have a distinguished point in $\sM_X^{\ch\theta'}$ degenerate
at $v$: the image under $\act_v$ of $t^{\ch\theta'} \in \msf Y^{\ch\theta',\ch\theta'}$. 
We will show that $\barM_X^{\ch\theta}$ contains this point. 
Then, stability of $\barM_X^{\ch\theta}$ under generic-Hecke correspondences 
(Theorem~\ref{thm:comp}(i) and Lemma~\ref{lem:act-equivariant}) 
and Proposition~\ref{prop:genHecke-transitive} imply that 
$\barM_X^{\ch\theta}$ contains all of $\sM_{X,v}^{\ch\theta'}$. 

\smallskip

Since $\ch\theta' \succeq \ch\theta$, we can decompose 
$\ch\theta' - \ch\theta = \sum_{j=1}^d \ch\nu_j$ for (not necessarily distinct) 
$\ch\nu_j$ equal to valuations of colors $D_j \in \Cal D$ 
(in case $D_j$ is not uniquely determined by its valuation, the choice of $D_j$ is arbitrary). 
The graded factorization property gives a map 
$\oo{\Scr C} := \sY^{\ch\theta,\ch\theta} \oo\xt \oo\prod_j \sY^{D_j}_{X^\bullet} 
\to \sY^{\ch\theta'}$. 
We claim that the image of 
$\oo{\Scr C}$ contains $C=\sY^{\ch\theta',\ch\theta'}_\red$ in its closure. 
(Note that $\oo{\Scr C}_\red = \oo C^{d+1}$.) 
This will produce an irreducible variety whose generic point maps to $\sM_X^{\ch\theta}$ 
while a special point maps to $\act_v(t^{\ch\theta'}) \in \sM_X^{\ch\theta'}$.

\smallskip

Consider the Beilinson--Drinfeld affine Grassmannian $\Gr_{B,C^{d+1}}$, 
whose fiber over $d+1$ pairwise distinct points $(v_0,\dotsc, v_d)\in \oo C^{d+1}$ 
is $\prod_{j=0}^d \Gr_B$ while the fiber over $v_0=\dotsc=v_d$ is $\Gr_B$. 
By Lemma~\ref{lem:Y=ScapGr}, we have an isomorphism $\msf Y^{\ch\lambda}\cong \Gr_B^{\ch\lambda}
\xt_{\msf L X/\msf L^+B}(\msf L^+X/\msf L^+B)$,
which depends on a fixed base point $x_0 \in X$. 
Now a reduction to $\mbb G_m \bs \GL_2$ (see proof of Lemma~\ref{lem:Ycolor}
and Example~\ref{eg:Yfiber})
shows that the central fiber $\msf Y^{D_j}_{X^\bullet} = \pt$ is contained in 
$\msf L^+ N \cdot t^{\ch\nu_j} \subset \Gr_B^{\nu_j}$.
Therefore, $\oo{\Scr C}$ is contained in 
the orbit of the multi-point jet space $(\Scr L^+N)_{C^{d+1}}$ acting on 
the closed subscheme $C^{d+1}\subset \Gr_{T, C^{d+1}}\subset \Gr_{B,C^{d+1}}$
given by $(v_0,\dotsc, v_d) \mapsto t_{v_0}^{\ch\theta} \prod_j t_{v_j}^{\ch\nu_j}$, where
the $v_j$'s are allowed to collide. 
The orbit of $(v_0,\dotsc, v_d) \in \oo C^{d+1}$ is 
\[ \textstyle
\{t_{v_0}^{\ch\theta}\} \xt \prod_{j=1}^d (\msf L^+ N \cdot t^{\ch\nu_j}_{v_j}) 
\subset \Gr_{B,v_0}^{\ch\theta} \xt \prod_{j=1}^d \Gr_{B,v_j}^{\ch\nu_j},
\]
while the orbit of the diagonal $v=v_0=\dotsc=v_d$ 
is $\msf L^+ N \cdot t^{\ch\theta'}_v = \{ t^{\ch\theta'}_v \} \subset \Gr_{B,v}^{\ch\theta'}$
since $\ch\theta'$ is antidominant. 

Now assume that $C=\mbb A^1$, which is justified by Proposition~\ref{prop:changecurve}.
Then we can identify $(\Scr L^+ N)_C = \msf L^+ N \times C$. 
Let $\msf Y^{D_j}_{X^\bullet}$ correspond to the point $n_j t^{\ch\nu_j}\in \Gr_B$
for $n_j \in \msf L^+N(k)$. 
For any pairwise distinct $v_0,\dotsc, v_d \in \mbb A^1$ we have
a line $\mbb A^1 \to C^{d+1} : a \mapsto (a v_0, \dotsc, a v_d)$ contracting all points to $0$. 
Multiplication defines a map 
$m: (\msf L^+ N \times C)^d \to (\Scr L^+ N)_{C^d}$. 
Letting $m(n_1,\dotsc,n_d)$ act 
on the point $t_{a v_0}^{\ch\theta} \prod_j t_{a v_j}^{\ch\nu_j} 
\in \Gr_{T,C^{d+1}}$ as $a\to 0$, 
we get a curve connecting $\{ t^{\ch\theta}_{v_0} \} \xt \prod \msf Y^{D_j}_{X^\bullet,v_j}$ to 
$t^{\ch\theta'}_0$. 
Hence the closure of $\oo{\Scr C}$ 
in $\sY^{\ch\theta'} \xt_{\sA^{\ch\theta'}} C^{d+1} \subset \Gr_{B,C^{d+1}}$ contains $\sY^{\ch\theta',\ch\theta'}_\red$.
Since the map $C^{d+1} \to \sA^{\ch\theta'}$ is finite and $\oo{\Scr C}$ is irreducible,
we have proved the claim.
\end{proof}

Now we draw some corollaries from Proposition \ref{prop:closure-rel}.


\begin{cor}
\source{intersection-complexes-unramified-l-factors}
\notready  \label{cor:can-comp}
\source{intersection-complexes-unramified-l-factors}
The open substack $\sM_X^0 = \Bun_H$ is dense in $\Scr M_{X^\can}$ iff $\mathfrak c_X^{\Cal D} \cap \Cal V = \mf c_{X^\can}^-$. 
\end{cor}

This is an analog of \cite[Proposition 1.2.3]{BG}, which says that 
$\Bun_B$ is dense in $\ol\Bun_B$ if $[G,G]$ is simply connected. 

\begin{proof}
Indeed, $\mathfrak c_X^{\Cal D} \cap \Cal V = \mf c_{X^\can}^-$ is precisely the condition that every $\ch\Theta \in \Sym^\infty(\mathfrak c_X^- \sm 0)$ is $\succeq 0$.
\end{proof}


Define $\Cal D^G_{\mathrm{sat}}(X)$ to be 
the set of primitive elements in $\Prim(\mathfrak c_X^-)$ that
cannot be decomposed as a sum $\ch\theta+\ch\lambda$ where 
$\ch\theta,\ch\lambda$ are both nonzero, 
$\ch\theta \in \mathfrak c_X^-$
and $\ch\lambda \succeq 0$ (see \S\ref{def:partition} for 
the definition of \emph{primitive}). 
Note that $\Cal D^G_{\mathrm{sat}}(X)$ contains $\varrho_X(\Cal D(X)\sm \Cal D)$, 
the valuations of the $G$-stable prime divisors,
but the containment may be strict.



\begin{cor}
\source{intersection-complexes-unramified-l-factors}
\notready \label{cor:comp}
\source{intersection-complexes-unramified-l-factors}
 There is a bijection between the set of irreducible
components of $\Scr M_X$ and 
\[ \pi_1(H) \xt \Sym^\infty(\Cal D^G_{\mathrm{sat}}(X)), \] 
such that $\ch\mu \in \pi_1(H),\, \ch\Theta \in  \Sym^\infty(\Cal D^G_{\mathrm{sat}}(X))$ corresponds to 
$^{\ch\mu}\barM_X^{\ch\Theta}$.  
\end{cor}

\begin{proof}
 For any $\ch\Theta' \in \Sym^\infty(\mathfrak c_X^- \sm 0)$, 
let $\ch\Theta'' \in \Sym^\infty(\mathfrak c_X^-\sm 0)$ 
be a minimal element with respect to the ordering $\preceq$
such that $\ch\Theta'' \preceq \ch\Theta'$. 
Then $\ch\Theta''$ can be refined to an element $\ch\Theta\in \Sym^\infty(\Cal D^G_{\mathrm{sat}}(X))$. Therefore, the closure relations
from Proposition~\ref{prop:closure-rel} tell us that any stratum 
is contained in the closure of $\sM^{\ch\Theta}_X$ for a partition 
$\ch\Theta$ as above. 
By definition if $\ch\theta_1,\ch\theta_2\in \Cal D^G_{\mathrm{sat}}(X)$
satisfy $\ch\theta_1\succeq \ch\theta_2$ then they must be equal. 
From this one deduces that $\sM^{\ch\Theta}_X$ is not contained
in the closure of any other stratum. 

Thus, the closure of each
$\sM^{\ch\Theta}_X,\, \ch\Theta\in \Sym^\infty(\Cal D^G_{\mathrm{sat}}(X))$, 
is a union of irreducible components, and no irreducible component 
is contained in two different such closures. The Corollary now follows from the description of irreducible components of $\sM^{\ch\Theta}_X$ by Corollary \ref{cor:strata-comp}.
\end{proof}


\begin{lem}
\source{intersection-complexes-unramified-l-factors}
\notready\label{lem:Membedding}
\source{intersection-complexes-unramified-l-factors}
Let $X_1 \to X_2$ be a $G$-equivariant map of affine spherical varieties with $X_1^\bullet = X_2^\bullet = H\bs G$. Then the induced map $\Scr M_{X_1} \to \Scr M_{X_2}$ is a closed embedding.
\end{lem}
\begin{proof}
Let $S$ be a test scheme and let $(\Scr P_G,\sigma_2) \in \Scr M_{X_2}(S)$, where
$\sigma_2 : C\xt S \to X_2 \xt^G \Scr P_G$ is a section. 
By Corollary~\ref{cor:Mcover} there exists a $B$-structure $\Scr P_B$ 
on $\Scr P_G$ after a suitable surjective \'etale base change $S'\to S$
such that $(\Scr P_B, \sigma_2)\in \Scr Y_{X_2}(S')$. Then in particular
there exists a relative effective divisor $D \subset C\xt S'$ such that 
$\sigma_2(C\xt S' \sm D) \subset X_2^\circ \xt^B \Scr P_B$, cf.~\S\ref{sss:YGr}. 
Since $X_1^\bullet=X_2^\bullet$, we have $X_1^\circ=X_2^\circ$. 
By Lemma~\ref{lem:extendsection} the condition that $\sigma_2|_{C\xt S'\sm D}$
extends to a section $C \xt S' \to X_1 \xt^B \Scr P_B = X_1 \xt^G \Scr P_G$ is
closed in $S'$. Therefore, $S' \xt_{\Scr M_{X_2}} \Scr M_{X_1} \to S'$ is a closed
embedding.
\end{proof}

Lemma~\ref{lem:Membedding} implies that 
$\sM_{X^\can}$ is a closed substack of $\sM_X$ containing $\barM_X^0$, and when the condition
of Corollary~\ref{cor:can-comp} is satisfied, $\sM_{X^\can} = \barM_X^0$.  


\subsection{Closure relations and components in the Zastava model}

In order to extend the results above to strata of $\sY$, we will need the following result:


\begin{lem}
\source{intersection-complexes-unramified-l-factors}
\notready \label{lem:Yraise}
\source{intersection-complexes-unramified-l-factors}
Let $y \in \sY^{\ch\lambda}(k)$ for $\ch\lambda \in \mf c_X$.
For any simple root $\alpha$ with $\Cal D(\alpha) = \{ D_\alpha^+, D_\alpha^- \}$ and any $N \gg 0$, there exists
a $k$-point 
\[ y' \in \sY^{\ch\lambda} \oo\xt C^{(N)} \oo\xt C^{(N)} \subset 
\sY^{\ch\lambda} \oo\xt \sY^{N\ch\nu_\alpha^+} \oo\xt \sY^{N\ch\nu_\alpha^-} 
\] 
such that the first coordinate is $y$ and the image of $y'$ under 
the composition
\[ 
    \sY^{\ch\lambda} \oo\xt \sY^{N\ch\nu_\alpha^+} \oo \xt
    \sY^{N\ch\nu_\alpha^-} \to 
    \sY^{\ch\lambda+N\ch\alpha} \to \sM_X 
\]
coincides with the image of $y$.
\end{lem}

Above we are using the fact, from Lemma~\ref{lem:Ycolor}, that 
$\sY^{N\ch\nu_\alpha^\pm}$ contains $C^{(N)}$ if $\ch\nu_\alpha^+ \ne \ch\nu_\alpha^-$ 
or $C^{(N)}\sqcup C^{(N)}$ if $\ch\nu_\alpha^+ = \ch\nu_\alpha^-$. 
In the latter case, Lemma~\ref{lem:Yraise} is picking out one of these components
(depending on the point $y$). 

\begin{proof}
The point $y$ is equivalent to a datum $(\Scr P_B, \sigma)$
where $\Scr P_B$ is a $B$-bundle on $C$ and $\sigma : C \to X \xt^B \Scr P_B$ is a section. 
Let $\Scr P_G = G \xt^B \Scr P_B$ denote the induced $G$-bundle.
The image of $y$ in $\sM_X$ corresponds to $(\Scr P_G, \sigma)$.
Set $\Bbbk := k(C)$. Then $\sigma|_{\Spec \Bbbk}$ defines a trivialization 
of $\Scr P_G|_{\Spec \Bbbk}$, which we fix. 
With respect to this trivialization, $B$-structures on $\Scr P_G$
are in bijection with sections $\Spec \Bbbk \to G/B$,
and $\Scr P_B$ corresponds to $1B \in (G/B)(\Bbbk)$. 
The preimage of $(\Scr P_G,\sigma)$ in $\sY$ identifies with the
orbit 
$H(\Bbbk)\cdot 1B \subset (G/B)(\Bbbk)$. 
Since $X^\circ P_\alpha/\mf R(P_\alpha)=\mbb G_m\bs \PGL_2$, the orbit of 
$H\cap P_\alpha$ on $1\in P_\alpha/B = \mbb P^1$ is 
$\mbb G_m$. Let $r \in \mbb G_m(\Bbbk) = k(C)^\times$ 
be a rational function on the curve $C$. 
Then the principal divisor defined by $r$ is of the form 
$\underline v^+ - \underline v^-$ where $\underline v^\pm$ 
are effective divisors with disjoint supports.
By the Riemann--Roch theorem, for any $N \gg 0$  
there exists $r \in k(C)^\times$ such that $\deg(\underline v^+)=\deg(\underline v^-)=N$ and the supports of $\underline v^+$ and $\underline v^-$
are both contained in $\sigma^{-1}(X^\circ \xt^B \Scr P_B)$.
Under the isomorphism $X^\circ P_\alpha/B \cong \mbb G_m \bs \mbb P^1$,
let us identify $D_\alpha^+$ with $0\in \mbb P^1$ and
$D_\alpha^-$ with $\infty \in \mbb P^1$. 
Then the preimage of $(\Scr P_G,\sigma)$ in $\sY$ corresponding to 
$r\in \mbb P^1(\Bbbk)$ has the desired property.
\end{proof}


Recall from Corollary \ref{cor:strata-comp} that the irreducible components of $\sM_X^{\ch\Theta}$ are denoted by ${^{\ch\mu}\sM_X^{\ch\Theta}}$, with $\ch\mu\in \pi_1(H)$. 
For any $\ch\lambda\in \mathfrak c_X$, define 
$^{\ch\mu}\sY^{\ch\lambda,\ch\Theta} = \sY^{\ch\lambda} \xt_{\sM_X} {^{\ch\mu}\sM_X^{\ch\Theta}}$
and 
$^{\ch\mu}\sY^{\ch\lambda,\succeq \ch\Theta} := \sY^{\ch\lambda} \xt_{\sM_X} 
{^{\ch\mu}}\barM_X^{\ch\Theta}$.


\begin{cor}
\source{intersection-complexes-unramified-l-factors}
\notready \label{cor:Ycomp}
\source{intersection-complexes-unramified-l-factors}
For any $\ch\lambda \in \mathfrak c_X,\, \ch\mu\in \pi_1(H),\,
\ch\Theta \in \Sym^\infty(\mathfrak c_X^- \sm 0)$, the scheme 
$^{\ch\mu}\sY^{\ch\lambda,\succeq\ch\Theta}$ is irreducible (if nonempty), and $^{\ch\mu}\sY^{\ch\lambda,\ch\Theta} $ is open dense in it.
\end{cor}
\begin{proof}
For $\ch\lambda$ large enough, the claim immediately follows from
Corollaries~\ref{cor:Mcover} and \ref{cor:comp}. 
Now, for arbitrary $\ch\lambda$, consider 
$\ch\lambda'=\ch\lambda+\sum_{\alpha\in \Delta_G} n_\alpha \ch\alpha$
large enough. 
By the graded factorization property and Lemma~\ref{lem:Yraise}, 
there exists an \'etale map 
\begin{equation} \label{e:Ycomp1}
\source{intersection-complexes-unramified-l-factors}
    ^{\ch\mu}\sY^{\ch\lambda,\succeq\ch\Theta} \oo\xt \Scr C \to {^{\ch\mu}}\sY^{\ch\lambda',\succeq\ch\Theta},
 \end{equation}
where $\Scr C = \oo\prod_{\Delta_G} C^{(n_\alpha)} \oo\xt C^{(n_\alpha)}$.


By the validity of the proposition for large $\ch\lambda'$,  
we know that $^{\ch\mu}\sY^{\ch\lambda',\ch\Theta}$ is 
connected and dense in $^{\ch\mu}\sY^{\ch\lambda', \succeq \ch\Theta}$.
Therefore, if $^{\ch\mu}\sY^{\ch\lambda,\succeq\ch\Theta}$ is not irreducible, 
there exist $y_1,y_2$ in disjoint connected components of 
$^{\ch\mu}\sY^{\ch\lambda,\ch\Theta}$ 
and $c_1,c_2\in \Scr C$ such that $(y_1,c_1)$ and $(y_2,c_2)$ have the same
image in $^{\ch\mu}\sY^{\ch\lambda',\ch\Theta}$.
Let $\abs{c_1}, \abs{c_2} \subset \abs C$ denote the support of $c_1,c_2$
as divisors. 
From the definition of the factorization map \eqref{e:Ycomp1}, 
if $(y_1,c_1)$ and $(y_2,c_2)$ have the same image, 
we have 
\[ y' = y_1 |_{C \sm (\abs{c_1}\cup \abs{c_2})} = y_2 |_{C\sm (\abs{c_1}\cup \abs{c_2})} : C \sm (\abs{c_1}\cup \abs{c_2}) \to X/B. \] 
We can extend $y'$ to an element of $\sY^{\ch\lambda-\ch\nu, \ch\Theta}$
with $\abs{c_1}\cup \abs{c_2}$ in its $B$-nondegenerate locus, 
for some $\ch\nu\succ 0$. 
By replacing $\Scr C$ by a possibly different irreducible scheme $\Scr C'$,
we may assume that $\abs{c_1}$ and $\abs{c_2}$ are disjoint. 
Now we must have $(y', c_2) \mapsto y_1,\, (y',c_1)\mapsto y_2$ under
the factorization map 
$\sY^{\ch\lambda-\ch\nu,\ch\Theta} \oo\xt \Scr C' \to \sY^{\ch\lambda,\ch\Theta}$.
This implies that $y_1,y_2$ were originally in the same connected component.
\end{proof}



\begin{cor}
\source{intersection-complexes-unramified-l-factors}
\notready\label{cor:compY}
\source{intersection-complexes-unramified-l-factors}
For every $\ch\lambda\in \mathfrak c_X$, there is an injection from the set of irreducible components of $\sY^{\ch\lambda}$ to $\pi_1(H) \xt \Sym^\infty(\Cal D^G_{\mathrm{sat}}(X))$.
\end{cor}

\begin{proof}
By Corollaries \ref{cor:Ycomp} and \ref{cor:comp}, the irreducible components $\sY^{\ch\lambda}$ are those $^{\ch\mu}\sY^{\ch\lambda,\succeq\ch\Theta}$ that are nonempty, for $\ch\Theta \in \Sym^\infty(\Cal D^G_{\mathrm{sat}}(X))$.
\end{proof}

In \S \ref{sect:connected-open}, and in particular Corollary \ref{cor:Ycomp2}, we will see a different description of the irreducible components of $\sY_X$, based on the partition of $\sY_{X^\bullet} = \sY_X^{?,0}$ into the subschemes $\sY^D_{X^\bullet}$ of \S \ref{sect:Y0cc}.
